\section{Covariance bounds for physically masked wave fields}
\label{S:masked-covariance}

\begin{figure}[htbp]
  \centering
  % Smooth toy profile and eta-squared mask; no covariance assertion.
% The unsmoothed potential is a distribution, not the curve drawn here.
% Reproduce with private/audit/2026-11-15-mask-window-drawing-held/generate.py.
\begin{tikzpicture}[x=2cm,y=1cm,font=\normalsize]
  \path[use as bounding box] (-2.85,-.25) rectangle (3.05,5.75);
  \foreach \xx in {-2,-1,1,2} {
  \draw[hamblue!22,densely dotted,line width=.35pt]
  (\xx,.66)--(\xx,5.06);
  }
  \foreach \yy in {4.62,2.50,1.05} {
  \draw[hamblue!24,line width=.3pt] (-2.65,\yy)--(2.65,\yy);
  }
  \node[anchor=west,inner sep=0pt] at (-2.65,5.52)
  {potential (after smoothing)};
  \node[anchor=west,inner sep=0pt] at (-2.65,3.66) {mask};
  \node[anchor=west,inner sep=0pt] at (-2.65,1.93) {what the wave sees};
  % potential: linear samples of the smooth toy function.
  \draw[hamblue,line width=.85pt,line join=round] plot coordinates {
  (-2.650000000,4.699833428)
  (-2.640000000,4.715627481)
  (-2.630000000,4.729790029)
  (-2.620000000,4.742266002)
  (-2.610000000,4.753026412)
  (-2.600000000,4.762067922)
  (-2.590000000,4.769411946)
  (-2.580000000,4.775103303)
  (-2.570000000,4.779208458)
  (-2.560000000,4.781813397)
  (-2.550000000,4.783021174)
  (-2.540000000,4.782949197)
  (-2.530000000,4.781726305)
  (-2.520000000,4.779489701)
  (-2.510000000,4.776381802)
  (-2.500000000,4.772547086)
  (-2.490000000,4.768128987)
  (-2.480000000,4.763266905)
  (-2.470000000,4.758093402)
  (-2.460000000,4.752731629)
  (-2.450000000,4.747293037)
  (-2.440000000,4.741875427)
  (-2.430000000,4.736561358)
  (-2.420000000,4.731416957)
  (-2.410000000,4.726491143)
  (-2.400000000,4.721815279)
  (-2.390000000,4.717403250)
  (-2.380000000,4.713251967)
  (-2.370000000,4.709342264)
  (-2.360000000,4.705640187)
  (-2.350000000,4.702098611)
  (-2.340000000,4.698659173)
  (-2.330000000,4.695254455)
  (-2.320000000,4.691810367)
  (-2.310000000,4.688248681)
  (-2.300000000,4.684489643)
  (-2.290000000,4.680454617)
  (-2.280000000,4.676068684)
  (-2.270000000,4.671263144)
  (-2.260000000,4.665977854)
  (-2.250000000,4.660163366)
  (-2.240000000,4.653782783)
  (-2.230000000,4.646813322)
  (-2.220000000,4.639247523)
  (-2.210000000,4.631094083)
  (-2.200000000,4.622378297)
  (-2.190000000,4.613142080)
  (-2.180000000,4.603443588)
  (-2.170000000,4.593356419)
  (-2.160000000,4.582968429)
  (-2.150000000,4.572380171)
  (-2.140000000,4.561703003)
  (-2.130000000,4.551056890)
  (-2.120000000,4.540567963)
  (-2.110000000,4.530365866)
  (-2.100000000,4.520580971)
  (-2.090000000,4.511341504)
  (-2.080000000,4.502770654)
  (-2.070000000,4.494983723)
  (-2.060000000,4.488085389)
  (-2.050000000,4.482167132)
  (-2.040000000,4.477304891)
  (-2.030000000,4.473557006)
  (-2.020000000,4.470962485)
  (-2.010000000,4.469539647)
  (-2.000000000,4.469285171)
  (-1.990000000,4.470173585)
  (-1.980000000,4.472157196)
  (-1.970000000,4.475166492)
  (-1.960000000,4.479110991)
  (-1.950000000,4.483880545)
  (-1.940000000,4.489347069)
  (-1.930000000,4.495366666)
  (-1.920000000,4.501782117)
  (-1.910000000,4.508425676)
  (-1.900000000,4.515122140)
  (-1.890000000,4.521692108)
  (-1.880000000,4.527955390)
  (-1.870000000,4.533734488)
  (-1.860000000,4.538858080)
  (-1.850000000,4.543164443)
  (-1.840000000,4.546504742)
  (-1.830000000,4.548746120)
  (-1.820000000,4.549774530)
  (-1.810000000,4.549497242)
  (-1.800000000,4.547844980)
  (-1.790000000,4.544773649)
  (-1.780000000,4.540265594)
  (-1.770000000,4.534330392)
  (-1.760000000,4.527005129)
  (-1.750000000,4.518354175)
  (-1.740000000,4.508468445)
  (-1.730000000,4.497464161)
  (-1.720000000,4.485481141)
  (-1.710000000,4.472680625)
  (-1.700000000,4.459242709)
  (-1.690000000,4.445363402)
  (-1.680000000,4.431251384)
  (-1.670000000,4.417124507)
  (-1.660000000,4.403206119)
  (-1.650000000,4.389721276)
  (-1.640000000,4.376892897)
  (-1.630000000,4.364937969)
  (-1.620000000,4.354063834)
  (-1.610000000,4.344464653)
  (-1.600000000,4.336318108)
  (-1.590000000,4.329782399)
  (-1.580000000,4.324993592)
  (-1.570000000,4.322063374)
  (-1.560000000,4.321077243)
  (-1.550000000,4.322093177)
  (-1.540000000,4.325140800)
  (-1.530000000,4.330221056)
  (-1.520000000,4.337306397)
  (-1.510000000,4.346341473)
  (-1.500000000,4.357244315)
  (-1.490000000,4.369907979)
  (-1.480000000,4.384202615)
  (-1.470000000,4.399977920)
  (-1.460000000,4.417065927)
  (-1.450000000,4.435284066)
  (-1.440000000,4.454438442)
  (-1.430000000,4.474327258)
  (-1.420000000,4.494744323)
  (-1.410000000,4.515482569)
  (-1.400000000,4.536337513)
  (-1.390000000,4.557110591)
  (-1.380000000,4.577612309)
  (-1.370000000,4.597665146)
  (-1.360000000,4.617106155)
  (-1.350000000,4.635789218)
  (-1.340000000,4.653586914)
  (-1.330000000,4.670391971)
  (-1.320000000,4.686118269)
  (-1.310000000,4.700701406)
  (-1.300000000,4.714098787)
  (-1.290000000,4.726289285)
  (-1.280000000,4.737272448)
  (-1.270000000,4.747067308)
  (-1.260000000,4.755710805)
  (-1.250000000,4.763255882)
  (-1.240000000,4.769769284)
  (-1.230000000,4.775329123)
  (-1.220000000,4.780022260)
  (-1.210000000,4.783941579)
  (-1.200000000,4.787183192)
  (-1.190000000,4.789843660)
  (-1.180000000,4.792017286)
  (-1.170000000,4.793793531)
  (-1.160000000,4.795254628)
  (-1.150000000,4.796473431)
  (-1.140000000,4.797511557)
  (-1.130000000,4.798417849)
  (-1.120000000,4.799227213)
  (-1.110000000,4.799959830)
  (-1.100000000,4.800620779)
  (-1.090000000,4.801200062)
  (-1.080000000,4.801673041)
  (-1.070000000,4.802001271)
  (-1.060000000,4.802133705)
  (-1.050000000,4.802008250)
  (-1.040000000,4.801553629)
  (-1.030000000,4.800691512)
  (-1.020000000,4.799338869)
  (-1.010000000,4.797410476)
  (-1.000000000,4.794821531)
  (-0.990000000,4.791490319)
  (-0.980000000,4.787340847)
  (-0.970000000,4.782305407)
  (-0.960000000,4.776326996)
  (-0.950000000,4.769361533)
  (-0.940000000,4.761379829)
  (-0.930000000,4.752369252)
  (-0.920000000,4.742335050)
  (-0.910000000,4.731301292)
  (-0.900000000,4.719311407)
  (-0.890000000,4.706428297)
  (-0.880000000,4.692734014)
  (-0.870000000,4.678329006)
  (-0.860000000,4.663330933)
  (-0.850000000,4.647873088)
  (-0.840000000,4.632102421)
  (-0.830000000,4.616177237)
  (-0.820000000,4.600264586)
  (-0.810000000,4.584537400)
  (-0.800000000,4.569171444)
  (-0.790000000,4.554342127)
  (-0.780000000,4.540221251)
  (-0.770000000,4.526973748)
  (-0.760000000,4.514754490)
  (-0.750000000,4.503705222)
  (-0.740000000,4.493951695)
  (-0.730000000,4.485601050)
  (-0.720000000,4.478739516)
  (-0.710000000,4.473430469)
  (-0.700000000,4.469712895)
  (-0.690000000,4.467600301)
  (-0.680000000,4.467080079)
  (-0.670000000,4.468113375)
  (-0.660000000,4.470635440)
  (-0.650000000,4.474556472)
  (-0.640000000,4.479762951)
  (-0.630000000,4.486119423)
  (-0.620000000,4.493470716)
  (-0.610000000,4.501644553)
  (-0.600000000,4.510454497)
  (-0.590000000,4.519703197)
  (-0.580000000,4.529185857)
  (-0.570000000,4.538693866)
  (-0.560000000,4.548018531)
  (-0.550000000,4.556954829)
  (-0.540000000,4.565305118)
  (-0.530000000,4.572882722)
  (-0.520000000,4.579515337)
  (-0.510000000,4.585048187)
  (-0.500000000,4.589346860)
  (-0.490000000,4.592299783)
  (-0.480000000,4.593820284)
  (-0.470000000,4.593848197)
  (-0.460000000,4.592350985)
  (-0.450000000,4.589324357)
  (-0.440000000,4.584792366)
  (-0.430000000,4.578806993)
  (-0.420000000,4.571447215)
  (-0.410000000,4.562817590)
  (-0.400000000,4.553046367)
  (-0.390000000,4.542283185)
  (-0.380000000,4.530696375)
  (-0.370000000,4.518469947)
  (-0.360000000,4.505800298)
  (-0.350000000,4.492892719)
  (-0.340000000,4.479957762)
  (-0.330000000,4.467207536)
  (-0.320000000,4.454852011)
  (-0.310000000,4.443095389)
  (-0.300000000,4.432132615)
  (-0.290000000,4.422146107)
  (-0.280000000,4.413302735)
  (-0.270000000,4.405751139)
  (-0.260000000,4.399619409)
  (-0.250000000,4.395013183)
  (-0.240000000,4.392014183)
  (-0.230000000,4.390679224)
  (-0.220000000,4.391039701)
  (-0.210000000,4.393101562)
  (-0.200000000,4.396845761)
  (-0.190000000,4.402229171)
  (-0.180000000,4.409185938)
  (-0.170000000,4.417629242)
  (-0.160000000,4.427453418)
  (-0.150000000,4.438536396)
  (-0.140000000,4.450742400)
  (-0.130000000,4.463924854)
  (-0.120000000,4.477929418)
  (-0.110000000,4.492597112)
  (-0.100000000,4.507767437)
  (-0.090000000,4.523281449)
  (-0.080000000,4.538984705)
  (-0.070000000,4.554730037)
  (-0.060000000,4.570380083)
  (-0.050000000,4.585809539)
  (-0.040000000,4.600907076)
  (-0.030000000,4.615576896)
  (-0.020000000,4.629739886)
  (-0.010000000,4.643334371)
  (0.000000000,4.656316433)
  (0.010000000,4.668659812)
  (0.020000000,4.680355393)
  (0.030000000,4.691410284)
  (0.040000000,4.701846541)
  (0.050000000,4.711699533)
  (0.060000000,4.721016032)
  (0.070000000,4.729852038)
  (0.080000000,4.738270414)
  (0.090000000,4.746338386)
  (0.100000000,4.754124960)
  (0.110000000,4.761698323)
  (0.120000000,4.769123288)
  (0.130000000,4.776458846)
  (0.140000000,4.783755885)
  (0.150000000,4.791055120)
  (0.160000000,4.798385298)
  (0.170000000,4.805761716)
  (0.180000000,4.813185077)
  (0.190000000,4.820640742)
  (0.200000000,4.828098365)
  (0.210000000,4.835511948)
  (0.220000000,4.842820307)
  (0.230000000,4.849947946)
  (0.240000000,4.856806320)
  (0.250000000,4.863295472)
  (0.260000000,4.869305995)
  (0.270000000,4.874721299)
  (0.280000000,4.879420116)
  (0.290000000,4.883279207)
  (0.300000000,4.886176201)
  (0.310000000,4.887992509)
  (0.320000000,4.888616251)
  (0.330000000,4.887945134)
  (0.340000000,4.885889209)
  (0.350000000,4.882373448)
  (0.360000000,4.877340091)
  (0.370000000,4.870750694)
  (0.380000000,4.862587839)
  (0.390000000,4.852856460)
  (0.400000000,4.841584754)
  (0.410000000,4.828824645)
  (0.420000000,4.814651786)
  (0.430000000,4.799165093)
  (0.440000000,4.782485815)
  (0.450000000,4.764756137)
  (0.460000000,4.746137361)
  (0.470000000,4.726807673)
  (0.480000000,4.706959539)
  (0.490000000,4.686796794)
  (0.500000000,4.666531446)
  (0.510000000,4.646380285)
  (0.520000000,4.626561340)
  (0.530000000,4.607290261)
  (0.540000000,4.588776688)
  (0.550000000,4.571220687)
  (0.560000000,4.554809318)
  (0.570000000,4.539713389)
  (0.580000000,4.526084488)
  (0.590000000,4.514052324)
  (0.600000000,4.503722444)
  (0.610000000,4.495174375)
  (0.620000000,4.488460212)
  (0.630000000,4.483603706)
  (0.640000000,4.480599846)
  (0.650000000,4.479414966)
  (0.660000000,4.479987362)
  (0.670000000,4.482228419)
  (0.680000000,4.486024220)
  (0.690000000,4.491237617)
  (0.700000000,4.497710712)
  (0.710000000,4.505267714)
  (0.720000000,4.513718119)
  (0.730000000,4.522860133)
  (0.740000000,4.532484304)
  (0.750000000,4.542377264)
  (0.760000000,4.552325536)
  (0.770000000,4.562119315)
  (0.780000000,4.571556157)
  (0.790000000,4.580444517)
  (0.800000000,4.588607052)
  (0.810000000,4.595883643)
  (0.820000000,4.602134072)
  (0.830000000,4.607240306)
  (0.840000000,4.611108349)
  (0.850000000,4.613669630)
  (0.860000000,4.614881904)
  (0.870000000,4.614729647)
  (0.880000000,4.613223955)
  (0.890000000,4.610401933)
  (0.900000000,4.606325616)
  (0.910000000,4.601080424)
  (0.920000000,4.594773202)
  (0.930000000,4.587529875)
  (0.940000000,4.579492785)
  (0.950000000,4.570817744)
  (0.960000000,4.561670885)
  (0.970000000,4.552225361)
  (0.980000000,4.542657969)
  (0.990000000,4.533145757)
  (1.000000000,4.523862695)
  (1.010000000,4.514976461)
  (1.020000000,4.506645423)
  (1.030000000,4.499015851)
  (1.040000000,4.492219439)
  (1.050000000,4.486371162)
  (1.060000000,4.481567521)
  (1.070000000,4.477885191)
  (1.080000000,4.475380118)
  (1.090000000,4.474087046)
  (1.100000000,4.474019504)
  (1.110000000,4.475170231)
  (1.120000000,4.477512033)
  (1.130000000,4.480999031)
  (1.140000000,4.485568291)
  (1.150000000,4.491141773)
  (1.160000000,4.497628564)
  (1.170000000,4.504927338)
  (1.180000000,4.512928993)
  (1.190000000,4.521519385)
  (1.200000000,4.530582123)
  (1.210000000,4.540001344)
  (1.220000000,4.549664403)
  (1.230000000,4.559464439)
  (1.240000000,4.569302734)
  (1.250000000,4.579090831)
  (1.260000000,4.588752364)
  (1.270000000,4.598224539)
  (1.280000000,4.607459272)
  (1.290000000,4.616423918)
  (1.300000000,4.625101610)
  (1.310000000,4.633491178)
  (1.320000000,4.641606673)
  (1.330000000,4.649476479)
  (1.340000000,4.657142070)
  (1.350000000,4.664656409)
  (1.360000000,4.672082043)
  (1.370000000,4.679488939)
  (1.380000000,4.686952101)
  (1.390000000,4.694549035)
  (1.400000000,4.702357110)
  (1.410000000,4.710450880)
  (1.420000000,4.718899440)
  (1.430000000,4.727763856)
  (1.440000000,4.737094750)
  (1.450000000,4.746930091)
  (1.460000000,4.757293243)
  (1.470000000,4.768191318)
  (1.480000000,4.779613886)
  (1.490000000,4.791532062)
  (1.500000000,4.803898012)
  (1.510000000,4.816644881)
  (1.520000000,4.829687167)
  (1.530000000,4.842921527)
  (1.540000000,4.856228015)
  (1.550000000,4.869471729)
  (1.560000000,4.882504838)
  (1.570000000,4.895168955)
  (1.580000000,4.907297811)
  (1.590000000,4.918720177)
  (1.600000000,4.929262981)
  (1.610000000,4.938754557)
  (1.620000000,4.947027956)
  (1.630000000,4.953924269)
  (1.640000000,4.959295866)
  (1.650000000,4.963009517)
  (1.660000000,4.964949298)
  (1.670000000,4.965019253)
  (1.680000000,4.963145723)
  (1.690000000,4.959279317)
  (1.700000000,4.953396467)
  (1.710000000,4.945500534)
  (1.720000000,4.935622439)
  (1.730000000,4.923820802)
  (1.740000000,4.910181574)
  (1.750000000,4.894817174)
  (1.760000000,4.877865130)
  (1.770000000,4.859486259)
  (1.780000000,4.839862400)
  (1.790000000,4.819193757)
  (1.800000000,4.797695885)
  (1.810000000,4.775596379)
  (1.820000000,4.753131330)
  (1.830000000,4.730541606)
  (1.840000000,4.708069028)
  (1.850000000,4.685952524)
  (1.860000000,4.664424311)
  (1.870000000,4.643706200)
  (1.880000000,4.624006072)
  (1.890000000,4.605514608)
  (1.900000000,4.588402322)
  (1.910000000,4.572816960)
  (1.920000000,4.558881308)
  (1.930000000,4.546691451)
  (1.940000000,4.536315519)
  (1.950000000,4.527792928)
  (1.960000000,4.521134149)
  (1.970000000,4.516320987)
  (1.980000000,4.513307376)
  (1.990000000,4.512020665)
  (2.000000000,4.512363369)
  (2.010000000,4.514215352)
  (2.020000000,4.517436392)
  (2.030000000,4.521869081)
  (2.040000000,4.527342001)
  (2.050000000,4.533673108)
  (2.060000000,4.540673266)
  (2.070000000,4.548149852)
  (2.080000000,4.555910374)
  (2.090000000,4.563766016)
  (2.100000000,4.571535055)
  (2.110000000,4.579046083)
  (2.120000000,4.586140967)
  (2.130000000,4.592677498)
  (2.140000000,4.598531677)
  (2.150000000,4.603599608)
  (2.160000000,4.607798943)
  (2.170000000,4.611069886)
  (2.180000000,4.613375718)
  (2.190000000,4.614702850)
  (2.200000000,4.615060410)
  (2.210000000,4.614479375)
  (2.220000000,4.613011279)
  (2.230000000,4.610726518)
  (2.240000000,4.607712314)
  (2.250000000,4.604070356)
  (2.260000000,4.599914206)
  (2.270000000,4.595366496)
  (2.280000000,4.590556003)
  (2.290000000,4.585614652)
  (2.300000000,4.580674514)
  (2.310000000,4.575864871)
  (2.320000000,4.571309402)
  (2.330000000,4.567123547)
  (2.340000000,4.563412123)
  (2.350000000,4.560267211)
  (2.360000000,4.557766386)
  (2.370000000,4.555971306)
  (2.380000000,4.554926691)
  (2.390000000,4.554659724)
  (2.400000000,4.555179854)
  (2.410000000,4.556479040)
  (2.420000000,4.558532382)
  (2.430000000,4.561299167)
  (2.440000000,4.564724262)
  (2.450000000,4.568739849)
  (2.460000000,4.573267446)
  (2.470000000,4.578220163)
  (2.480000000,4.583505153)
  (2.490000000,4.589026186)
  (2.500000000,4.594686300)
  (2.510000000,4.600390448)
  (2.520000000,4.606048103)
  (2.530000000,4.611575736)
  (2.540000000,4.616899132)
  (2.550000000,4.621955472)
  (2.560000000,4.626695133)
  (2.570000000,4.631083179)
  (2.580000000,4.635100487)
  (2.590000000,4.638744487)
  (2.600000000,4.642029498)
  (2.610000000,4.644986650)
  (2.620000000,4.647663379)
  (2.630000000,4.650122517)
  (2.640000000,4.652440981)
  (2.650000000,4.654708096)
  };
  % mask: linear samples of the smooth toy function.
  \draw[hamteal,line width=.85pt,line join=round] plot coordinates {
  (-2.650000000,2.500000000)
  (-2.640000000,2.500000000)
  (-2.630000000,2.500000000)
  (-2.620000000,2.500000000)
  (-2.610000000,2.500000000)
  (-2.600000000,2.500000000)
  (-2.590000000,2.500000000)
  (-2.580000000,2.500000000)
  (-2.570000000,2.500000000)
  (-2.560000000,2.500000000)
  (-2.550000000,2.500000000)
  (-2.540000000,2.500000000)
  (-2.530000000,2.500000000)
  (-2.520000000,2.500000000)
  (-2.510000000,2.500000000)
  (-2.500000000,2.500000000)
  (-2.490000000,2.500000000)
  (-2.480000000,2.500000000)
  (-2.470000000,2.500000000)
  (-2.460000000,2.500000000)
  (-2.450000000,2.500000000)
  (-2.440000000,2.500000000)
  (-2.430000000,2.500000000)
  (-2.420000000,2.500000000)
  (-2.410000000,2.500000000)
  (-2.400000000,2.500000000)
  (-2.390000000,2.500000000)
  (-2.380000000,2.500000000)
  (-2.370000000,2.500000000)
  (-2.360000000,2.500000000)
  (-2.350000000,2.500000000)
  (-2.340000000,2.500000000)
  (-2.330000000,2.500000000)
  (-2.320000000,2.500000000)
  (-2.310000000,2.500000000)
  (-2.300000000,2.500000000)
  (-2.290000000,2.500000000)
  (-2.280000000,2.500000000)
  (-2.270000000,2.500000000)
  (-2.260000000,2.500000000)
  (-2.250000000,2.500000000)
  (-2.240000000,2.500000000)
  (-2.230000000,2.500000000)
  (-2.220000000,2.500000000)
  (-2.210000000,2.500000000)
  (-2.200000000,2.500000000)
  (-2.190000000,2.500000000)
  (-2.180000000,2.500000000)
  (-2.170000000,2.500000000)
  (-2.160000000,2.500000000)
  (-2.150000000,2.500000000)
  (-2.140000000,2.500000000)
  (-2.130000000,2.500000000)
  (-2.120000000,2.500000000)
  (-2.110000000,2.500000000)
  (-2.100000000,2.500000000)
  (-2.090000000,2.500000000)
  (-2.080000000,2.500000000)
  (-2.070000000,2.500000000)
  (-2.060000000,2.500000000)
  (-2.050000000,2.500000000)
  (-2.040000000,2.500000000)
  (-2.030000000,2.500000000)
  (-2.020000000,2.500000000)
  (-2.010000000,2.500000000)
  (-2.000000000,2.500000000)
  (-1.990000000,2.500000000)
  (-1.980000000,2.500000000)
  (-1.970000000,2.500000000)
  (-1.960000000,2.500000000)
  (-1.950000000,2.500000000)
  (-1.940000000,2.500000000)
  (-1.930000000,2.500000000)
  (-1.920000000,2.500000000)
  (-1.910000000,2.500000001)
  (-1.900000000,2.500000012)
  (-1.890000000,2.500000078)
  (-1.880000000,2.500000364)
  (-1.870000000,2.500001345)
  (-1.860000000,2.500004135)
  (-1.850000000,2.500010980)
  (-1.840000000,2.500025868)
  (-1.830000000,2.500055210)
  (-1.820000000,2.500108487)
  (-1.810000000,2.500198793)
  (-1.800000000,2.500343174)
  (-1.790000000,2.500562748)
  (-1.780000000,2.500882565)
  (-1.770000000,2.501331226)
  (-1.760000000,2.501940310)
  (-1.750000000,2.502743645)
  (-1.740000000,2.503776499)
  (-1.730000000,2.505074730)
  (-1.720000000,2.506673974)
  (-1.710000000,2.508608893)
  (-1.700000000,2.510912529)
  (-1.690000000,2.513615779)
  (-1.680000000,2.516746996)
  (-1.670000000,2.520331720)
  (-1.660000000,2.524392532)
  (-1.650000000,2.528949012)
  (-1.640000000,2.534017780)
  (-1.630000000,2.539612618)
  (-1.620000000,2.545744635)
  (-1.610000000,2.552422468)
  (-1.600000000,2.559652500)
  (-1.590000000,2.567439081)
  (-1.580000000,2.575784736)
  (-1.570000000,2.584690363)
  (-1.560000000,2.594155389)
  (-1.550000000,2.604177904)
  (-1.540000000,2.614754752)
  (-1.530000000,2.625881573)
  (-1.520000000,2.637552811)
  (-1.510000000,2.649761666)
  (-1.500000000,2.662500000)
  (-1.490000000,2.675758195)
  (-1.480000000,2.689524962)
  (-1.470000000,2.703787097)
  (-1.460000000,2.718529199)
  (-1.450000000,2.733733333)
  (-1.440000000,2.749378660)
  (-1.430000000,2.765441022)
  (-1.420000000,2.781892503)
  (-1.410000000,2.798700965)
  (-1.400000000,2.815829569)
  (-1.390000000,2.833236313)
  (-1.380000000,2.850873571)
  (-1.370000000,2.868687695)
  (-1.360000000,2.886618676)
  (-1.350000000,2.904599909)
  (-1.340000000,2.922558095)
  (-1.330000000,2.940413319)
  (-1.320000000,2.958079344)
  (-1.310000000,2.975464175)
  (-1.300000000,2.992470919)
  (-1.290000000,3.008998995)
  (-1.280000000,3.024945721)
  (-1.270000000,3.040208290)
  (-1.260000000,3.054686146)
  (-1.250000000,3.068283726)
  (-1.240000000,3.080913520)
  (-1.230000000,3.092499370)
  (-1.220000000,3.102979859)
  (-1.210000000,3.112311657)
  (-1.200000000,3.120472593)
  (-1.190000000,3.127464222)
  (-1.180000000,3.133313628)
  (-1.170000000,3.138074151)
  (-1.160000000,3.141824777)
  (-1.150000000,3.144667886)
  (-1.140000000,3.146725171)
  (-1.130000000,3.148131607)
  (-1.120000000,3.149027585)
  (-1.110000000,3.149549643)
  (-1.100000000,3.149820750)
  (-1.090000000,3.149941700)
  (-1.080000000,3.149985635)
  (-1.070000000,3.149997619)
  (-1.060000000,3.149999782)
  (-1.050000000,3.149999992)
  (-1.040000000,3.150000000)
  (-1.030000000,3.150000000)
  (-1.020000000,3.150000000)
  (-1.010000000,3.150000000)
  (-1.000000000,3.150000000)
  (-0.990000000,3.150000000)
  (-0.980000000,3.150000000)
  (-0.970000000,3.150000000)
  (-0.960000000,3.150000000)
  (-0.950000000,3.150000000)
  (-0.940000000,3.150000000)
  (-0.930000000,3.150000000)
  (-0.920000000,3.150000000)
  (-0.910000000,3.150000000)
  (-0.900000000,3.150000000)
  (-0.890000000,3.150000000)
  (-0.880000000,3.150000000)
  (-0.870000000,3.150000000)
  (-0.860000000,3.150000000)
  (-0.850000000,3.150000000)
  (-0.840000000,3.150000000)
  (-0.830000000,3.150000000)
  (-0.820000000,3.150000000)
  (-0.810000000,3.150000000)
  (-0.800000000,3.150000000)
  (-0.790000000,3.150000000)
  (-0.780000000,3.150000000)
  (-0.770000000,3.150000000)
  (-0.760000000,3.150000000)
  (-0.750000000,3.150000000)
  (-0.740000000,3.150000000)
  (-0.730000000,3.150000000)
  (-0.720000000,3.150000000)
  (-0.710000000,3.150000000)
  (-0.700000000,3.150000000)
  (-0.690000000,3.150000000)
  (-0.680000000,3.150000000)
  (-0.670000000,3.150000000)
  (-0.660000000,3.150000000)
  (-0.650000000,3.150000000)
  (-0.640000000,3.150000000)
  (-0.630000000,3.150000000)
  (-0.620000000,3.150000000)
  (-0.610000000,3.150000000)
  (-0.600000000,3.150000000)
  (-0.590000000,3.150000000)
  (-0.580000000,3.150000000)
  (-0.570000000,3.150000000)
  (-0.560000000,3.150000000)
  (-0.550000000,3.150000000)
  (-0.540000000,3.150000000)
  (-0.530000000,3.150000000)
  (-0.520000000,3.150000000)
  (-0.510000000,3.150000000)
  (-0.500000000,3.150000000)
  (-0.490000000,3.150000000)
  (-0.480000000,3.150000000)
  (-0.470000000,3.150000000)
  (-0.460000000,3.150000000)
  (-0.450000000,3.150000000)
  (-0.440000000,3.150000000)
  (-0.430000000,3.150000000)
  (-0.420000000,3.150000000)
  (-0.410000000,3.150000000)
  (-0.400000000,3.150000000)
  (-0.390000000,3.150000000)
  (-0.380000000,3.150000000)
  (-0.370000000,3.150000000)
  (-0.360000000,3.150000000)
  (-0.350000000,3.150000000)
  (-0.340000000,3.150000000)
  (-0.330000000,3.150000000)
  (-0.320000000,3.150000000)
  (-0.310000000,3.150000000)
  (-0.300000000,3.150000000)
  (-0.290000000,3.150000000)
  (-0.280000000,3.150000000)
  (-0.270000000,3.150000000)
  (-0.260000000,3.150000000)
  (-0.250000000,3.150000000)
  (-0.240000000,3.150000000)
  (-0.230000000,3.150000000)
  (-0.220000000,3.150000000)
  (-0.210000000,3.150000000)
  (-0.200000000,3.150000000)
  (-0.190000000,3.150000000)
  (-0.180000000,3.150000000)
  (-0.170000000,3.150000000)
  (-0.160000000,3.150000000)
  (-0.150000000,3.150000000)
  (-0.140000000,3.150000000)
  (-0.130000000,3.150000000)
  (-0.120000000,3.150000000)
  (-0.110000000,3.150000000)
  (-0.100000000,3.150000000)
  (-0.090000000,3.150000000)
  (-0.080000000,3.150000000)
  (-0.070000000,3.150000000)
  (-0.060000000,3.150000000)
  (-0.050000000,3.150000000)
  (-0.040000000,3.150000000)
  (-0.030000000,3.150000000)
  (-0.020000000,3.150000000)
  (-0.010000000,3.150000000)
  (0.000000000,3.150000000)
  (0.010000000,3.150000000)
  (0.020000000,3.150000000)
  (0.030000000,3.150000000)
  (0.040000000,3.150000000)
  (0.050000000,3.150000000)
  (0.060000000,3.150000000)
  (0.070000000,3.150000000)
  (0.080000000,3.150000000)
  (0.090000000,3.150000000)
  (0.100000000,3.150000000)
  (0.110000000,3.150000000)
  (0.120000000,3.150000000)
  (0.130000000,3.150000000)
  (0.140000000,3.150000000)
  (0.150000000,3.150000000)
  (0.160000000,3.150000000)
  (0.170000000,3.150000000)
  (0.180000000,3.150000000)
  (0.190000000,3.150000000)
  (0.200000000,3.150000000)
  (0.210000000,3.150000000)
  (0.220000000,3.150000000)
  (0.230000000,3.150000000)
  (0.240000000,3.150000000)
  (0.250000000,3.150000000)
  (0.260000000,3.150000000)
  (0.270000000,3.150000000)
  (0.280000000,3.150000000)
  (0.290000000,3.150000000)
  (0.300000000,3.150000000)
  (0.310000000,3.150000000)
  (0.320000000,3.150000000)
  (0.330000000,3.150000000)
  (0.340000000,3.150000000)
  (0.350000000,3.150000000)
  (0.360000000,3.150000000)
  (0.370000000,3.150000000)
  (0.380000000,3.150000000)
  (0.390000000,3.150000000)
  (0.400000000,3.150000000)
  (0.410000000,3.150000000)
  (0.420000000,3.150000000)
  (0.430000000,3.150000000)
  (0.440000000,3.150000000)
  (0.450000000,3.150000000)
  (0.460000000,3.150000000)
  (0.470000000,3.150000000)
  (0.480000000,3.150000000)
  (0.490000000,3.150000000)
  (0.500000000,3.150000000)
  (0.510000000,3.150000000)
  (0.520000000,3.150000000)
  (0.530000000,3.150000000)
  (0.540000000,3.150000000)
  (0.550000000,3.150000000)
  (0.560000000,3.150000000)
  (0.570000000,3.150000000)
  (0.580000000,3.150000000)
  (0.590000000,3.150000000)
  (0.600000000,3.150000000)
  (0.610000000,3.150000000)
  (0.620000000,3.150000000)
  (0.630000000,3.150000000)
  (0.640000000,3.150000000)
  (0.650000000,3.150000000)
  (0.660000000,3.150000000)
  (0.670000000,3.150000000)
  (0.680000000,3.150000000)
  (0.690000000,3.150000000)
  (0.700000000,3.150000000)
  (0.710000000,3.150000000)
  (0.720000000,3.150000000)
  (0.730000000,3.150000000)
  (0.740000000,3.150000000)
  (0.750000000,3.150000000)
  (0.760000000,3.150000000)
  (0.770000000,3.150000000)
  (0.780000000,3.150000000)
  (0.790000000,3.150000000)
  (0.800000000,3.150000000)
  (0.810000000,3.150000000)
  (0.820000000,3.150000000)
  (0.830000000,3.150000000)
  (0.840000000,3.150000000)
  (0.850000000,3.150000000)
  (0.860000000,3.150000000)
  (0.870000000,3.150000000)
  (0.880000000,3.150000000)
  (0.890000000,3.150000000)
  (0.900000000,3.150000000)
  (0.910000000,3.150000000)
  (0.920000000,3.150000000)
  (0.930000000,3.150000000)
  (0.940000000,3.150000000)
  (0.950000000,3.150000000)
  (0.960000000,3.150000000)
  (0.970000000,3.150000000)
  (0.980000000,3.150000000)
  (0.990000000,3.150000000)
  (1.000000000,3.150000000)
  (1.010000000,3.150000000)
  (1.020000000,3.150000000)
  (1.030000000,3.150000000)
  (1.040000000,3.150000000)
  (1.050000000,3.149999992)
  (1.060000000,3.149999782)
  (1.070000000,3.149997619)
  (1.080000000,3.149985635)
  (1.090000000,3.149941700)
  (1.100000000,3.149820750)
  (1.110000000,3.149549643)
  (1.120000000,3.149027585)
  (1.130000000,3.148131607)
  (1.140000000,3.146725171)
  (1.150000000,3.144667886)
  (1.160000000,3.141824777)
  (1.170000000,3.138074151)
  (1.180000000,3.133313628)
  (1.190000000,3.127464222)
  (1.200000000,3.120472593)
  (1.210000000,3.112311657)
  (1.220000000,3.102979859)
  (1.230000000,3.092499370)
  (1.240000000,3.080913520)
  (1.250000000,3.068283726)
  (1.260000000,3.054686146)
  (1.270000000,3.040208290)
  (1.280000000,3.024945721)
  (1.290000000,3.008998995)
  (1.300000000,2.992470919)
  (1.310000000,2.975464175)
  (1.320000000,2.958079344)
  (1.330000000,2.940413319)
  (1.340000000,2.922558095)
  (1.350000000,2.904599909)
  (1.360000000,2.886618676)
  (1.370000000,2.868687695)
  (1.380000000,2.850873571)
  (1.390000000,2.833236313)
  (1.400000000,2.815829569)
  (1.410000000,2.798700965)
  (1.420000000,2.781892503)
  (1.430000000,2.765441022)
  (1.440000000,2.749378660)
  (1.450000000,2.733733333)
  (1.460000000,2.718529199)
  (1.470000000,2.703787097)
  (1.480000000,2.689524962)
  (1.490000000,2.675758195)
  (1.500000000,2.662500000)
  (1.510000000,2.649761666)
  (1.520000000,2.637552811)
  (1.530000000,2.625881573)
  (1.540000000,2.614754752)
  (1.550000000,2.604177904)
  (1.560000000,2.594155389)
  (1.570000000,2.584690363)
  (1.580000000,2.575784736)
  (1.590000000,2.567439081)
  (1.600000000,2.559652500)
  (1.610000000,2.552422468)
  (1.620000000,2.545744635)
  (1.630000000,2.539612618)
  (1.640000000,2.534017780)
  (1.650000000,2.528949012)
  (1.660000000,2.524392532)
  (1.670000000,2.520331720)
  (1.680000000,2.516746996)
  (1.690000000,2.513615779)
  (1.700000000,2.510912529)
  (1.710000000,2.508608893)
  (1.720000000,2.506673974)
  (1.730000000,2.505074730)
  (1.740000000,2.503776499)
  (1.750000000,2.502743645)
  (1.760000000,2.501940310)
  (1.770000000,2.501331226)
  (1.780000000,2.500882565)
  (1.790000000,2.500562748)
  (1.800000000,2.500343174)
  (1.810000000,2.500198793)
  (1.820000000,2.500108487)
  (1.830000000,2.500055210)
  (1.840000000,2.500025868)
  (1.850000000,2.500010980)
  (1.860000000,2.500004135)
  (1.870000000,2.500001345)
  (1.880000000,2.500000364)
  (1.890000000,2.500000078)
  (1.900000000,2.500000012)
  (1.910000000,2.500000001)
  (1.920000000,2.500000000)
  (1.930000000,2.500000000)
  (1.940000000,2.500000000)
  (1.950000000,2.500000000)
  (1.960000000,2.500000000)
  (1.970000000,2.500000000)
  (1.980000000,2.500000000)
  (1.990000000,2.500000000)
  (2.000000000,2.500000000)
  (2.010000000,2.500000000)
  (2.020000000,2.500000000)
  (2.030000000,2.500000000)
  (2.040000000,2.500000000)
  (2.050000000,2.500000000)
  (2.060000000,2.500000000)
  (2.070000000,2.500000000)
  (2.080000000,2.500000000)
  (2.090000000,2.500000000)
  (2.100000000,2.500000000)
  (2.110000000,2.500000000)
  (2.120000000,2.500000000)
  (2.130000000,2.500000000)
  (2.140000000,2.500000000)
  (2.150000000,2.500000000)
  (2.160000000,2.500000000)
  (2.170000000,2.500000000)
  (2.180000000,2.500000000)
  (2.190000000,2.500000000)
  (2.200000000,2.500000000)
  (2.210000000,2.500000000)
  (2.220000000,2.500000000)
  (2.230000000,2.500000000)
  (2.240000000,2.500000000)
  (2.250000000,2.500000000)
  (2.260000000,2.500000000)
  (2.270000000,2.500000000)
  (2.280000000,2.500000000)
  (2.290000000,2.500000000)
  (2.300000000,2.500000000)
  (2.310000000,2.500000000)
  (2.320000000,2.500000000)
  (2.330000000,2.500000000)
  (2.340000000,2.500000000)
  (2.350000000,2.500000000)
  (2.360000000,2.500000000)
  (2.370000000,2.500000000)
  (2.380000000,2.500000000)
  (2.390000000,2.500000000)
  (2.400000000,2.500000000)
  (2.410000000,2.500000000)
  (2.420000000,2.500000000)
  (2.430000000,2.500000000)
  (2.440000000,2.500000000)
  (2.450000000,2.500000000)
  (2.460000000,2.500000000)
  (2.470000000,2.500000000)
  (2.480000000,2.500000000)
  (2.490000000,2.500000000)
  (2.500000000,2.500000000)
  (2.510000000,2.500000000)
  (2.520000000,2.500000000)
  (2.530000000,2.500000000)
  (2.540000000,2.500000000)
  (2.550000000,2.500000000)
  (2.560000000,2.500000000)
  (2.570000000,2.500000000)
  (2.580000000,2.500000000)
  (2.590000000,2.500000000)
  (2.600000000,2.500000000)
  (2.610000000,2.500000000)
  (2.620000000,2.500000000)
  (2.630000000,2.500000000)
  (2.640000000,2.500000000)
  (2.650000000,2.500000000)
  };
  % Product of each pair of displayed linear segments, as cubic Bezier paths.
  \draw[hamblue,line width=.85pt,line join=round] (-2.650000000,1.050000000)
  .. controls (-2.646666667,1.050000000)
  and (-2.643333333,1.050000000)
  .. (-2.640000000,1.050000000)
  .. controls (-2.636666667,1.050000000)
  and (-2.633333333,1.050000000)
  .. (-2.630000000,1.050000000)
  .. controls (-2.626666667,1.050000000)
  and (-2.623333333,1.050000000)
  .. (-2.620000000,1.050000000)
  .. controls (-2.616666667,1.050000000)
  and (-2.613333333,1.050000000)
  .. (-2.610000000,1.050000000)
  .. controls (-2.606666667,1.050000000)
  and (-2.603333333,1.050000000)
  .. (-2.600000000,1.050000000)
  .. controls (-2.596666667,1.050000000)
  and (-2.593333333,1.050000000)
  .. (-2.590000000,1.050000000)
  .. controls (-2.586666667,1.050000000)
  and (-2.583333333,1.050000000)
  .. (-2.580000000,1.050000000)
  .. controls (-2.576666667,1.050000000)
  and (-2.573333333,1.050000000)
  .. (-2.570000000,1.050000000)
  .. controls (-2.566666667,1.050000000)
  and (-2.563333333,1.050000000)
  .. (-2.560000000,1.050000000)
  .. controls (-2.556666667,1.050000000)
  and (-2.553333333,1.050000000)
  .. (-2.550000000,1.050000000)
  .. controls (-2.546666667,1.050000000)
  and (-2.543333333,1.050000000)
  .. (-2.540000000,1.050000000)
  .. controls (-2.536666667,1.050000000)
  and (-2.533333333,1.050000000)
  .. (-2.530000000,1.050000000)
  .. controls (-2.526666667,1.050000000)
  and (-2.523333333,1.050000000)
  .. (-2.520000000,1.050000000)
  .. controls (-2.516666667,1.050000000)
  and (-2.513333333,1.050000000)
  .. (-2.510000000,1.050000000)
  .. controls (-2.506666667,1.050000000)
  and (-2.503333333,1.050000000)
  .. (-2.500000000,1.050000000)
  .. controls (-2.496666667,1.050000000)
  and (-2.493333333,1.050000000)
  .. (-2.490000000,1.050000000)
  .. controls (-2.486666667,1.050000000)
  and (-2.483333333,1.050000000)
  .. (-2.480000000,1.050000000)
  .. controls (-2.476666667,1.050000000)
  and (-2.473333333,1.050000000)
  .. (-2.470000000,1.050000000)
  .. controls (-2.466666667,1.050000000)
  and (-2.463333333,1.050000000)
  .. (-2.460000000,1.050000000)
  .. controls (-2.456666667,1.050000000)
  and (-2.453333333,1.050000000)
  .. (-2.450000000,1.050000000)
  .. controls (-2.446666667,1.050000000)
  and (-2.443333333,1.050000000)
  .. (-2.440000000,1.050000000)
  .. controls (-2.436666667,1.050000000)
  and (-2.433333333,1.050000000)
  .. (-2.430000000,1.050000000)
  .. controls (-2.426666667,1.050000000)
  and (-2.423333333,1.050000000)
  .. (-2.420000000,1.050000000)
  .. controls (-2.416666667,1.050000000)
  and (-2.413333333,1.050000000)
  .. (-2.410000000,1.050000000)
  .. controls (-2.406666667,1.050000000)
  and (-2.403333333,1.050000000)
  .. (-2.400000000,1.050000000)
  .. controls (-2.396666667,1.050000000)
  and (-2.393333333,1.050000000)
  .. (-2.390000000,1.050000000)
  .. controls (-2.386666667,1.050000000)
  and (-2.383333333,1.050000000)
  .. (-2.380000000,1.050000000)
  .. controls (-2.376666667,1.050000000)
  and (-2.373333333,1.050000000)
  .. (-2.370000000,1.050000000)
  .. controls (-2.366666667,1.050000000)
  and (-2.363333333,1.050000000)
  .. (-2.360000000,1.050000000)
  .. controls (-2.356666667,1.050000000)
  and (-2.353333333,1.050000000)
  .. (-2.350000000,1.050000000)
  .. controls (-2.346666667,1.050000000)
  and (-2.343333333,1.050000000)
  .. (-2.340000000,1.050000000)
  .. controls (-2.336666667,1.050000000)
  and (-2.333333333,1.050000000)
  .. (-2.330000000,1.050000000)
  .. controls (-2.326666667,1.050000000)
  and (-2.323333333,1.050000000)
  .. (-2.320000000,1.050000000)
  .. controls (-2.316666667,1.050000000)
  and (-2.313333333,1.050000000)
  .. (-2.310000000,1.050000000)
  .. controls (-2.306666667,1.050000000)
  and (-2.303333333,1.050000000)
  .. (-2.300000000,1.050000000)
  .. controls (-2.296666667,1.050000000)
  and (-2.293333333,1.050000000)
  .. (-2.290000000,1.050000000)
  .. controls (-2.286666667,1.050000000)
  and (-2.283333333,1.050000000)
  .. (-2.280000000,1.050000000)
  .. controls (-2.276666667,1.050000000)
  and (-2.273333333,1.050000000)
  .. (-2.270000000,1.050000000)
  .. controls (-2.266666667,1.050000000)
  and (-2.263333333,1.050000000)
  .. (-2.260000000,1.050000000)
  .. controls (-2.256666667,1.050000000)
  and (-2.253333333,1.050000000)
  .. (-2.250000000,1.050000000)
  .. controls (-2.246666667,1.050000000)
  and (-2.243333333,1.050000000)
  .. (-2.240000000,1.050000000)
  .. controls (-2.236666667,1.050000000)
  and (-2.233333333,1.050000000)
  .. (-2.230000000,1.050000000)
  .. controls (-2.226666667,1.050000000)
  and (-2.223333333,1.050000000)
  .. (-2.220000000,1.050000000)
  .. controls (-2.216666667,1.050000000)
  and (-2.213333333,1.050000000)
  .. (-2.210000000,1.050000000)
  .. controls (-2.206666667,1.050000000)
  and (-2.203333333,1.050000000)
  .. (-2.200000000,1.050000000)
  .. controls (-2.196666667,1.050000000)
  and (-2.193333333,1.050000000)
  .. (-2.190000000,1.050000000)
  .. controls (-2.186666667,1.050000000)
  and (-2.183333333,1.050000000)
  .. (-2.180000000,1.050000000)
  .. controls (-2.176666667,1.050000000)
  and (-2.173333333,1.050000000)
  .. (-2.170000000,1.050000000)
  .. controls (-2.166666667,1.050000000)
  and (-2.163333333,1.050000000)
  .. (-2.160000000,1.050000000)
  .. controls (-2.156666667,1.050000000)
  and (-2.153333333,1.050000000)
  .. (-2.150000000,1.050000000)
  .. controls (-2.146666667,1.050000000)
  and (-2.143333333,1.050000000)
  .. (-2.140000000,1.050000000)
  .. controls (-2.136666667,1.050000000)
  and (-2.133333333,1.050000000)
  .. (-2.130000000,1.050000000)
  .. controls (-2.126666667,1.050000000)
  and (-2.123333333,1.050000000)
  .. (-2.120000000,1.050000000)
  .. controls (-2.116666667,1.050000000)
  and (-2.113333333,1.050000000)
  .. (-2.110000000,1.050000000)
  .. controls (-2.106666667,1.050000000)
  and (-2.103333333,1.050000000)
  .. (-2.100000000,1.050000000)
  .. controls (-2.096666667,1.050000000)
  and (-2.093333333,1.050000000)
  .. (-2.090000000,1.050000000)
  .. controls (-2.086666667,1.050000000)
  and (-2.083333333,1.050000000)
  .. (-2.080000000,1.050000000)
  .. controls (-2.076666667,1.050000000)
  and (-2.073333333,1.050000000)
  .. (-2.070000000,1.050000000)
  .. controls (-2.066666667,1.050000000)
  and (-2.063333333,1.050000000)
  .. (-2.060000000,1.050000000)
  .. controls (-2.056666667,1.050000000)
  and (-2.053333333,1.050000000)
  .. (-2.050000000,1.050000000)
  .. controls (-2.046666667,1.050000000)
  and (-2.043333333,1.050000000)
  .. (-2.040000000,1.050000000)
  .. controls (-2.036666667,1.050000000)
  and (-2.033333333,1.050000000)
  .. (-2.030000000,1.050000000)
  .. controls (-2.026666667,1.050000000)
  and (-2.023333333,1.050000000)
  .. (-2.020000000,1.050000000)
  .. controls (-2.016666667,1.050000000)
  and (-2.013333333,1.050000000)
  .. (-2.010000000,1.050000000)
  .. controls (-2.006666667,1.050000000)
  and (-2.003333333,1.050000000)
  .. (-2.000000000,1.050000000)
  .. controls (-1.996666667,1.050000000)
  and (-1.993333333,1.050000000)
  .. (-1.990000000,1.050000000)
  .. controls (-1.986666667,1.050000000)
  and (-1.983333333,1.050000000)
  .. (-1.980000000,1.050000000)
  .. controls (-1.976666667,1.050000000)
  and (-1.973333333,1.050000000)
  .. (-1.970000000,1.050000000)
  .. controls (-1.966666667,1.050000000)
  and (-1.963333333,1.050000000)
  .. (-1.960000000,1.050000000)
  .. controls (-1.956666667,1.050000000)
  and (-1.953333333,1.050000000)
  .. (-1.950000000,1.050000000)
  .. controls (-1.946666667,1.050000000)
  and (-1.943333333,1.050000000)
  .. (-1.940000000,1.050000000)
  .. controls (-1.936666667,1.050000000)
  and (-1.933333333,1.050000000)
  .. (-1.930000000,1.050000000)
  .. controls (-1.926666667,1.050000000)
  and (-1.923333333,1.050000000)
  .. (-1.920000000,1.050000000)
  .. controls (-1.916666667,1.050000000)
  and (-1.913333333,1.050000000)
  .. (-1.910000000,1.050000000)
  .. controls (-1.906666667,1.049999999)
  and (-1.903333333,1.049999999)
  .. (-1.900000000,1.049999998)
  .. controls (-1.896666667,1.049999995)
  and (-1.893333333,1.049999991)
  .. (-1.890000000,1.049999988)
  .. controls (-1.886666667,1.049999974)
  and (-1.883333333,1.049999961)
  .. (-1.880000000,1.049999948)
  .. controls (-1.876666667,1.049999903)
  and (-1.873333333,1.049999861)
  .. (-1.870000000,1.049999822)
  .. controls (-1.866666667,1.049999702)
  and (-1.863333333,1.049999589)
  .. (-1.860000000,1.049999484)
  .. controls (-1.856666667,1.049999208)
  and (-1.853333333,1.049998948)
  .. (-1.850000000,1.049998702)
  .. controls (-1.846666667,1.049998134)
  and (-1.843333333,1.049997592)
  .. (-1.840000000,1.049997075)
  .. controls (-1.836666667,1.049995999)
  and (-1.833333333,1.049994957)
  .. (-1.830000000,1.049993948)
  .. controls (-1.826666667,1.049992030)
  and (-1.823333333,1.049990141)
  .. (-1.820000000,1.049988279)
  .. controls (-1.816666667,1.049985012)
  and (-1.813333333,1.049981731)
  .. (-1.810000000,1.049978438)
  .. controls (-1.806666667,1.049973049)
  and (-1.803333333,1.049967538)
  .. (-1.800000000,1.049961905)
  .. controls (-1.796666667,1.049953240)
  and (-1.793333333,1.049944229)
  .. (-1.790000000,1.049934872)
  .. controls (-1.786666667,1.049921233)
  and (-1.783333333,1.049906855)
  .. (-1.780000000,1.049891737)
  .. controls (-1.776666667,1.049870705)
  and (-1.773333333,1.049848308)
  .. (-1.770000000,1.049824545)
  .. controls (-1.766666667,1.049792785)
  and (-1.763333333,1.049758738)
  .. (-1.760000000,1.049722402)
  .. controls (-1.756666667,1.049675483)
  and (-1.753333333,1.049625000)
  .. (-1.750000000,1.049570954)
  .. controls (-1.746666667,1.049503206)
  and (-1.743333333,1.049430222)
  .. (-1.740000000,1.049352002)
  .. controls (-1.736666667,1.049256437)
  and (-1.733333333,1.049153546)
  .. (-1.730000000,1.049043329)
  .. controls (-1.726666667,1.048911649)
  and (-1.723333333,1.048770142)
  .. (-1.720000000,1.048618807)
  .. controls (-1.716666667,1.048441518)
  and (-1.713333333,1.048251528)
  .. (-1.710000000,1.048048836)
  .. controls (-1.706666667,1.047815474)
  and (-1.703333333,1.047566237)
  .. (-1.700000000,1.047301125)
  .. controls (-1.696666667,1.047000599)
  and (-1.693333333,1.046680833)
  .. (-1.690000000,1.046341826)
  .. controls (-1.686666667,1.045962866)
  and (-1.683333333,1.045561246)
  .. (-1.680000000,1.045136966)
  .. controls (-1.676666667,1.044668661)
  and (-1.673333333,1.044174386)
  .. (-1.670000000,1.043654142)
  .. controls (-1.666666667,1.043086540)
  and (-1.663333333,1.042489953)
  .. (-1.660000000,1.041864382)
  .. controls (-1.656666667,1.041189127)
  and (-1.653333333,1.040482363)
  .. (-1.650000000,1.039744090)
  .. controls (-1.646666667,1.038955065)
  and (-1.643333333,1.038132695)
  .. (-1.640000000,1.037276979)
  .. controls (-1.636666667,1.036370915)
  and (-1.633333333,1.035430550)
  .. (-1.630000000,1.034455885)
  .. controls (-1.626666667,1.033432912)
  and (-1.623333333,1.032375744)
  .. (-1.620000000,1.031284380)
  .. controls (-1.616666667,1.030148489)
  and (-1.613333333,1.028979725)
  .. (-1.610000000,1.027778088)
  .. controls (-1.606666667,1.026537477)
  and (-1.603333333,1.025266661)
  .. (-1.600000000,1.023965640)
  .. controls (-1.596666667,1.022632930)
  and (-1.593333333,1.021274123)
  .. (-1.590000000,1.019889218)
  .. controls (-1.586666667,1.018481521)
  and (-1.583333333,1.017053329)
  .. (-1.580000000,1.015604642)
  .. controls (-1.576666667,1.014143471)
  and (-1.573333333,1.012668918)
  .. (-1.570000000,1.011180983)
  .. controls (-1.566666667,1.009692012)
  and (-1.563333333,1.008198254)
  .. (-1.560000000,1.006699710)
  .. controls (-1.556666667,1.005212376)
  and (-1.553333333,1.003730263)
  .. (-1.550000000,1.002253371)
  .. controls (-1.546666667,1.000800336)
  and (-1.543333333,0.999363830)
  .. (-1.540000000,0.997943855)
  .. controls (-1.536666667,0.996560336)
  and (-1.533333333,0.995205805)
  .. (-1.530000000,0.993880263)
  .. controls (-1.526666667,0.992603255)
  and (-1.523333333,0.991368655)
  .. (-1.520000000,0.990176462)
  .. controls (-1.516666667,0.989043865)
  and (-1.513333333,0.987967835)
  .. (-1.510000000,0.986948374)
  .. controls (-1.506666667,0.985998053)
  and (-1.503333333,0.985118955)
  .. (-1.500000000,0.984311079)
  .. controls (-1.496666667,0.983579889)
  and (-1.493333333,0.982934800)
  .. (-1.490000000,0.982375812)
  .. controls (-1.486666667,0.981898602)
  and (-1.483333333,0.981522311)
  .. (-1.480000000,0.981246938)
  .. controls (-1.476666667,0.981055574)
  and (-1.473333333,0.980979589)
  .. (-1.470000000,0.981018983)
  .. controls (-1.466666667,0.981141407)
  and (-1.463333333,0.981393018)
  .. (-1.460000000,0.981773815)
  .. controls (-1.456666667,0.982233178)
  and (-1.453333333,0.982834588)
  .. (-1.450000000,0.983578045)
  .. controls (-1.446666667,0.984391929)
  and (-1.443333333,0.985359494)
  .. (-1.440000000,0.986480739)
  .. controls (-1.436666667,0.987660502)
  and (-1.433333333,0.989004090)
  .. (-1.430000000,0.990511505)
  .. controls (-1.426666667,0.992061759)
  and (-1.423333333,0.993784264)
  .. (-1.420000000,0.995679021)
  .. controls (-1.416666667,0.997597278)
  and (-1.413333333,0.999694293)
  .. (-1.410000000,1.001970066)
  .. controls (-1.406666667,1.004246555)
  and (-1.403333333,1.006706232)
  .. (-1.400000000,1.009349097)
  .. controls (-1.396666667,1.011966769)
  and (-1.393333333,1.014769873)
  .. (-1.390000000,1.017758408)
  .. controls (-1.386666667,1.020693136)
  and (-1.383333333,1.023813298)
  .. (-1.380000000,1.027118892)
  .. controls (-1.376666667,1.030339872)
  and (-1.373333333,1.033744043)
  .. (-1.370000000,1.037331406)
  .. controls (-1.366666667,1.040801752)
  and (-1.363333333,1.044450866)
  .. (-1.360000000,1.048278747)
  .. controls (-1.356666667,1.051956278)
  and (-1.353333333,1.055806089)
  .. (-1.350000000,1.059828179)
  .. controls (-1.346666667,1.063666380)
  and (-1.343333333,1.067668486)
  .. (-1.340000000,1.071834496)
  .. controls (-1.336666667,1.075783632)
  and (-1.333333333,1.079886643)
  .. (-1.330000000,1.084143531)
  .. controls (-1.326666667,1.088151888)
  and (-1.323333333,1.092302718)
  .. (-1.320000000,1.096596021)
  .. controls (-1.316666667,1.100611246)
  and (-1.313333333,1.104756483)
  .. (-1.310000000,1.109031734)
  .. controls (-1.306666667,1.113002218)
  and (-1.303333333,1.117089546)
  .. (-1.300000000,1.121293717)
  .. controls (-1.296666667,1.125169993)
  and (-1.293333333,1.129149595)
  .. (-1.290000000,1.133232522)
  .. controls (-1.286666667,1.136968617)
  and (-1.283333333,1.140794530)
  .. (-1.280000000,1.144710262)
  .. controls (-1.276666667,1.148264953)
  and (-1.273333333,1.151896308)
  .. (-1.270000000,1.155604328)
  .. controls (-1.266666667,1.158942251)
  and (-1.263333333,1.162344349)
  .. (-1.260000000,1.165810620)
  .. controls (-1.256666667,1.168903178)
  and (-1.253333333,1.172048349)
  .. (-1.250000000,1.175246133)
  .. controls (-1.246666667,1.178072160)
  and (-1.243333333,1.180940373)
  .. (-1.240000000,1.183850772)
  .. controls (-1.236666667,1.186396921)
  and (-1.233333333,1.188976103)
  .. (-1.230000000,1.191588319)
  .. controls (-1.226666667,1.193849143)
  and (-1.223333333,1.196135190)
  .. (-1.220000000,1.198446461)
  .. controls (-1.216666667,1.200424187)
  and (-1.213333333,1.202420670)
  .. (-1.210000000,1.204435908)
  .. controls (-1.206666667,1.206139905)
  and (-1.203333333,1.207857468)
  .. (-1.200000000,1.209588598)
  .. controls (-1.196666667,1.211034562)
  and (-1.193333333,1.212490065)
  .. (-1.190000000,1.213955108)
  .. controls (-1.186666667,1.215164009)
  and (-1.183333333,1.216379430)
  .. (-1.180000000,1.217601371)
  .. controls (-1.176666667,1.218598198)
  and (-1.173333333,1.219599361)
  .. (-1.170000000,1.220604861)
  .. controls (-1.166666667,1.221417232)
  and (-1.163333333,1.222232413)
  .. (-1.160000000,1.223050404)
  .. controls (-1.156666667,1.223707084)
  and (-1.153333333,1.224365541)
  .. (-1.150000000,1.225025775)
  .. controls (-1.146666667,1.225555161)
  and (-1.143333333,1.226085642)
  .. (-1.140000000,1.226617218)
  .. controls (-1.136666667,1.227045824)
  and (-1.133333333,1.227475083)
  .. (-1.130000000,1.227904996)
  .. controls (-1.126666667,1.228255987)
  and (-1.123333333,1.228607350)
  .. (-1.120000000,1.228959085)
  .. controls (-1.116666667,1.229250908)
  and (-1.113333333,1.229542928)
  .. (-1.110000000,1.229835144)
  .. controls (-1.106666667,1.230080327)
  and (-1.103333333,1.230325603)
  .. (-1.100000000,1.230570970)
  .. controls (-1.096666667,1.230775214)
  and (-1.093333333,1.230979494)
  .. (-1.090000000,1.231183810)
  .. controls (-1.086666667,1.231345538)
  and (-1.083333333,1.231507277)
  .. (-1.080000000,1.231669026)
  .. controls (-1.076666667,1.231779550)
  and (-1.073333333,1.231890076)
  .. (-1.070000000,1.232000604)
  .. controls (-1.066666667,1.232044951)
  and (-1.063333333,1.232089298)
  .. (-1.060000000,1.232133644)
  .. controls (-1.056666667,1.232091846)
  and (-1.053333333,1.232050047)
  .. (-1.050000000,1.232008248)
  .. controls (-1.046666667,1.231856708)
  and (-1.043333333,1.231705168)
  .. (-1.040000000,1.231553629)
  .. controls (-1.036666667,1.231266257)
  and (-1.033333333,1.230978885)
  .. (-1.030000000,1.230691512)
  .. controls (-1.026666667,1.230240631)
  and (-1.023333333,1.229789750)
  .. (-1.020000000,1.229338869)
  .. controls (-1.016666667,1.228696071)
  and (-1.013333333,1.228053274)
  .. (-1.010000000,1.227410476)
  .. controls (-1.006666667,1.226547494)
  and (-1.003333333,1.225684512)
  .. (-1.000000000,1.224821531)
  .. controls (-0.996666667,1.223711127)
  and (-0.993333333,1.222600723)
  .. (-0.990000000,1.221490319)
  .. controls (-0.986666667,1.220107162)
  and (-0.983333333,1.218724004)
  .. (-0.980000000,1.217340847)
  .. controls (-0.976666667,1.215662367)
  and (-0.973333333,1.213983887)
  .. (-0.970000000,1.212305407)
  .. controls (-0.966666667,1.210312603)
  and (-0.963333333,1.208319800)
  .. (-0.960000000,1.206326996)
  .. controls (-0.956666667,1.204005175)
  and (-0.953333333,1.201683354)
  .. (-0.950000000,1.199361533)
  .. controls (-0.946666667,1.196700965)
  and (-0.943333333,1.194040397)
  .. (-0.940000000,1.191379829)
  .. controls (-0.936666667,1.188376303)
  and (-0.933333333,1.185372777)
  .. (-0.930000000,1.182369252)
  .. controls (-0.926666667,1.179024518)
  and (-0.923333333,1.175679784)
  .. (-0.920000000,1.172335050)
  .. controls (-0.916666667,1.168657130)
  and (-0.913333333,1.164979211)
  .. (-0.910000000,1.161301292)
  .. controls (-0.906666667,1.157304664)
  and (-0.903333333,1.153308035)
  .. (-0.900000000,1.149311407)
  .. controls (-0.896666667,1.145017037)
  and (-0.893333333,1.140722667)
  .. (-0.890000000,1.136428297)
  .. controls (-0.886666667,1.131863536)
  and (-0.883333333,1.127298775)
  .. (-0.880000000,1.122734014)
  .. controls (-0.876666667,1.117932345)
  and (-0.873333333,1.113130675)
  .. (-0.870000000,1.108329006)
  .. controls (-0.866666667,1.103329648)
  and (-0.863333333,1.098330291)
  .. (-0.860000000,1.093330933)
  .. controls (-0.856666667,1.088178318)
  and (-0.853333333,1.083025703)
  .. (-0.850000000,1.077873088)
  .. controls (-0.846666667,1.072616199)
  and (-0.843333333,1.067359310)
  .. (-0.840000000,1.062102421)
  .. controls (-0.836666667,1.056794026)
  and (-0.833333333,1.051485632)
  .. (-0.830000000,1.046177237)
  .. controls (-0.826666667,1.040873020)
  and (-0.823333333,1.035568803)
  .. (-0.820000000,1.030264586)
  .. controls (-0.816666667,1.025022191)
  and (-0.813333333,1.019779796)
  .. (-0.810000000,1.014537400)
  .. controls (-0.806666667,1.009415415)
  and (-0.803333333,1.004293429)
  .. (-0.800000000,0.999171444)
  .. controls (-0.796666667,0.994228338)
  and (-0.793333333,0.989285233)
  .. (-0.790000000,0.984342127)
  .. controls (-0.786666667,0.979635169)
  and (-0.783333333,0.974928210)
  .. (-0.780000000,0.970221251)
  .. controls (-0.776666667,0.965805417)
  and (-0.773333333,0.961389582)
  .. (-0.770000000,0.956973748)
  .. controls (-0.766666667,0.952900662)
  and (-0.763333333,0.948827576)
  .. (-0.760000000,0.944754490)
  .. controls (-0.756666667,0.941071400)
  and (-0.753333333,0.937388311)
  .. (-0.750000000,0.933705222)
  .. controls (-0.746666667,0.930454046)
  and (-0.743333333,0.927202870)
  .. (-0.740000000,0.923951695)
  .. controls (-0.736666667,0.921168146)
  and (-0.733333333,0.918384598)
  .. (-0.730000000,0.915601050)
  .. controls (-0.726666667,0.913313872)
  and (-0.723333333,0.911026694)
  .. (-0.720000000,0.908739516)
  .. controls (-0.716666667,0.906969833)
  and (-0.713333333,0.905200151)
  .. (-0.710000000,0.903430469)
  .. controls (-0.706666667,0.902191278)
  and (-0.703333333,0.900952087)
  .. (-0.700000000,0.899712895)
  .. controls (-0.696666667,0.899008697)
  and (-0.693333333,0.898304499)
  .. (-0.690000000,0.897600301)
  .. controls (-0.686666667,0.897426893)
  and (-0.683333333,0.897253486)
  .. (-0.680000000,0.897080079)
  .. controls (-0.676666667,0.897424511)
  and (-0.673333333,0.897768943)
  .. (-0.670000000,0.898113375)
  .. controls (-0.666666667,0.898954064)
  and (-0.663333333,0.899794752)
  .. (-0.660000000,0.900635440)
  .. controls (-0.656666667,0.901942451)
  and (-0.653333333,0.903249461)
  .. (-0.650000000,0.904556472)
  .. controls (-0.646666667,0.906291965)
  and (-0.643333333,0.908027458)
  .. (-0.640000000,0.909762951)
  .. controls (-0.636666667,0.911881775)
  and (-0.633333333,0.914000599)
  .. (-0.630000000,0.916119423)
  .. controls (-0.626666667,0.918569854)
  and (-0.623333333,0.921020285)
  .. (-0.620000000,0.923470716)
  .. controls (-0.616666667,0.926195329)
  and (-0.613333333,0.928919941)
  .. (-0.610000000,0.931644553)
  .. controls (-0.606666667,0.934581201)
  and (-0.603333333,0.937517849)
  .. (-0.600000000,0.940454497)
  .. controls (-0.596666667,0.943537397)
  and (-0.593333333,0.946620297)
  .. (-0.590000000,0.949703197)
  .. controls (-0.586666667,0.952864084)
  and (-0.583333333,0.956024970)
  .. (-0.580000000,0.959185857)
  .. controls (-0.576666667,0.962355193)
  and (-0.573333333,0.965524529)
  .. (-0.570000000,0.968693866)
  .. controls (-0.566666667,0.971802087)
  and (-0.563333333,0.974910309)
  .. (-0.560000000,0.978018531)
  .. controls (-0.556666667,0.980997297)
  and (-0.553333333,0.983976063)
  .. (-0.550000000,0.986954829)
  .. controls (-0.546666667,0.989738259)
  and (-0.543333333,0.992521688)
  .. (-0.540000000,0.995305118)
  .. controls (-0.536666667,0.997830986)
  and (-0.533333333,1.000356854)
  .. (-0.530000000,1.002882722)
  .. controls (-0.526666667,1.005093594)
  and (-0.523333333,1.007304465)
  .. (-0.520000000,1.009515337)
  .. controls (-0.516666667,1.011359621)
  and (-0.513333333,1.013203904)
  .. (-0.510000000,1.015048187)
  .. controls (-0.506666667,1.016481078)
  and (-0.503333333,1.017913969)
  .. (-0.500000000,1.019346860)
  .. controls (-0.496666667,1.020331168)
  and (-0.493333333,1.021315475)
  .. (-0.490000000,1.022299783)
  .. controls (-0.486666667,1.022806616)
  and (-0.483333333,1.023313450)
  .. (-0.480000000,1.023820284)
  .. controls (-0.476666667,1.023829588)
  and (-0.473333333,1.023838892)
  .. (-0.470000000,1.023848197)
  .. controls (-0.466666667,1.023349126)
  and (-0.463333333,1.022850055)
  .. (-0.460000000,1.022350985)
  .. controls (-0.456666667,1.021342109)
  and (-0.453333333,1.020333233)
  .. (-0.450000000,1.019324357)
  .. controls (-0.446666667,1.017813693)
  and (-0.443333333,1.016303029)
  .. (-0.440000000,1.014792366)
  .. controls (-0.436666667,1.012797241)
  and (-0.433333333,1.010802117)
  .. (-0.430000000,1.008806993)
  .. controls (-0.426666667,1.006353734)
  and (-0.423333333,1.003900474)
  .. (-0.420000000,1.001447215)
  .. controls (-0.416666667,0.998570673)
  and (-0.413333333,0.995694132)
  .. (-0.410000000,0.992817590)
  .. controls (-0.406666667,0.989560516)
  and (-0.403333333,0.986303442)
  .. (-0.400000000,0.983046367)
  .. controls (-0.396666667,0.979458640)
  and (-0.393333333,0.975870913)
  .. (-0.390000000,0.972283185)
  .. controls (-0.386666667,0.968420915)
  and (-0.383333333,0.964558645)
  .. (-0.380000000,0.960696375)
  .. controls (-0.376666667,0.956620899)
  and (-0.373333333,0.952545423)
  .. (-0.370000000,0.948469947)
  .. controls (-0.366666667,0.944246731)
  and (-0.363333333,0.940023514)
  .. (-0.360000000,0.935800298)
  .. controls (-0.356666667,0.931497772)
  and (-0.353333333,0.927195245)
  .. (-0.350000000,0.922892719)
  .. controls (-0.346666667,0.918581067)
  and (-0.343333333,0.914269414)
  .. (-0.340000000,0.909957762)
  .. controls (-0.336666667,0.905707687)
  and (-0.333333333,0.901457611)
  .. (-0.330000000,0.897207536)
  .. controls (-0.326666667,0.893089028)
  and (-0.323333333,0.888970520)
  .. (-0.320000000,0.884852011)
  .. controls (-0.316666667,0.880933137)
  and (-0.313333333,0.877014263)
  .. (-0.310000000,0.873095389)
  .. controls (-0.306666667,0.869441131)
  and (-0.303333333,0.865786873)
  .. (-0.300000000,0.862132615)
  .. controls (-0.296666667,0.858803779)
  and (-0.293333333,0.855474943)
  .. (-0.290000000,0.852146107)
  .. controls (-0.286666667,0.849198316)
  and (-0.283333333,0.846250526)
  .. (-0.280000000,0.843302735)
  .. controls (-0.276666667,0.840785536)
  and (-0.273333333,0.838268338)
  .. (-0.270000000,0.835751139)
  .. controls (-0.266666667,0.833707229)
  and (-0.263333333,0.831663319)
  .. (-0.260000000,0.829619409)
  .. controls (-0.256666667,0.828084001)
  and (-0.253333333,0.826548592)
  .. (-0.250000000,0.825013183)
  .. controls (-0.246666667,0.824013517)
  and (-0.243333333,0.823013850)
  .. (-0.240000000,0.822014183)
  .. controls (-0.236666667,0.821569197)
  and (-0.233333333,0.821124210)
  .. (-0.230000000,0.820679224)
  .. controls (-0.226666667,0.820799383)
  and (-0.223333333,0.820919542)
  .. (-0.220000000,0.821039701)
  .. controls (-0.216666667,0.821726988)
  and (-0.213333333,0.822414275)
  .. (-0.210000000,0.823101562)
  .. controls (-0.206666667,0.824349628)
  and (-0.203333333,0.825597695)
  .. (-0.200000000,0.826845761)
  .. controls (-0.196666667,0.828640231)
  and (-0.193333333,0.830434701)
  .. (-0.190000000,0.832229171)
  .. controls (-0.186666667,0.834548093)
  and (-0.183333333,0.836867016)
  .. (-0.180000000,0.839185938)
  .. controls (-0.176666667,0.842000373)
  and (-0.173333333,0.844814808)
  .. (-0.170000000,0.847629242)
  .. controls (-0.166666667,0.850903968)
  and (-0.163333333,0.854178693)
  .. (-0.160000000,0.857453418)
  .. controls (-0.156666667,0.861147744)
  and (-0.153333333,0.864842070)
  .. (-0.150000000,0.868536396)
  .. controls (-0.146666667,0.872605064)
  and (-0.143333333,0.876673732)
  .. (-0.140000000,0.880742400)
  .. controls (-0.136666667,0.885136551)
  and (-0.133333333,0.889530703)
  .. (-0.130000000,0.893924854)
  .. controls (-0.126666667,0.898593042)
  and (-0.123333333,0.903261230)
  .. (-0.120000000,0.907929418)
  .. controls (-0.116666667,0.912818649)
  and (-0.113333333,0.917707880)
  .. (-0.110000000,0.922597112)
  .. controls (-0.106666667,0.927653887)
  and (-0.103333333,0.932710662)
  .. (-0.100000000,0.937767437)
  .. controls (-0.096666667,0.942938774)
  and (-0.093333333,0.948110112)
  .. (-0.090000000,0.953281449)
  .. controls (-0.086666667,0.958515868)
  and (-0.083333333,0.963750287)
  .. (-0.080000000,0.968984705)
  .. controls (-0.076666667,0.974233149)
  and (-0.073333333,0.979481593)
  .. (-0.070000000,0.984730037)
  .. controls (-0.066666667,0.989946719)
  and (-0.063333333,0.995163401)
  .. (-0.060000000,1.000380083)
  .. controls (-0.056666667,1.005523235)
  and (-0.053333333,1.010666387)
  .. (-0.050000000,1.015809539)
  .. controls (-0.046666667,1.020842051)
  and (-0.043333333,1.025874564)
  .. (-0.040000000,1.030907076)
  .. controls (-0.036666667,1.035797016)
  and (-0.033333333,1.040686956)
  .. (-0.030000000,1.045576896)
  .. controls (-0.026666667,1.050297892)
  and (-0.023333333,1.055018889)
  .. (-0.020000000,1.059739886)
  .. controls (-0.016666667,1.064271381)
  and (-0.013333333,1.068802876)
  .. (-0.010000000,1.073334371)
  .. controls (-0.006666667,1.077661725)
  and (-0.003333333,1.081989079)
  .. (0.000000000,1.086316433)
  .. controls (0.003333333,1.090430893)
  and (0.006666667,1.094545353)
  .. (0.010000000,1.098659812)
  .. controls (0.013333333,1.102558339)
  and (0.016666667,1.106456866)
  .. (0.020000000,1.110355393)
  .. controls (0.023333333,1.114040357)
  and (0.026666667,1.117725321)
  .. (0.030000000,1.121410284)
  .. controls (0.033333333,1.124889036)
  and (0.036666667,1.128367789)
  .. (0.040000000,1.131846541)
  .. controls (0.043333333,1.135130871)
  and (0.046666667,1.138415202)
  .. (0.050000000,1.141699533)
  .. controls (0.053333333,1.144805033)
  and (0.056666667,1.147910533)
  .. (0.060000000,1.151016032)
  .. controls (0.063333333,1.153961368)
  and (0.066666667,1.156906703)
  .. (0.070000000,1.159852038)
  .. controls (0.073333333,1.162658163)
  and (0.076666667,1.165464288)
  .. (0.080000000,1.168270414)
  .. controls (0.083333333,1.170959738)
  and (0.086666667,1.173649062)
  .. (0.090000000,1.176338386)
  .. controls (0.093333333,1.178933911)
  and (0.096666667,1.181529436)
  .. (0.100000000,1.184124960)
  .. controls (0.103333333,1.186649415)
  and (0.106666667,1.189173869)
  .. (0.110000000,1.191698323)
  .. controls (0.113333333,1.194173311)
  and (0.116666667,1.196648300)
  .. (0.120000000,1.199123288)
  .. controls (0.123333333,1.201568474)
  and (0.126666667,1.204013660)
  .. (0.130000000,1.206458846)
  .. controls (0.133333333,1.208891192)
  and (0.136666667,1.211323538)
  .. (0.140000000,1.213755885)
  .. controls (0.143333333,1.216188963)
  and (0.146666667,1.218622041)
  .. (0.150000000,1.221055120)
  .. controls (0.153333333,1.223498513)
  and (0.156666667,1.225941906)
  .. (0.160000000,1.228385298)
  .. controls (0.163333333,1.230844104)
  and (0.166666667,1.233302910)
  .. (0.170000000,1.235761716)
  .. controls (0.173333333,1.238236169)
  and (0.176666667,1.240710623)
  .. (0.180000000,1.243185077)
  .. controls (0.183333333,1.245670298)
  and (0.186666667,1.248155520)
  .. (0.190000000,1.250640742)
  .. controls (0.193333333,1.253126616)
  and (0.196666667,1.255612490)
  .. (0.200000000,1.258098365)
  .. controls (0.203333333,1.260569559)
  and (0.206666667,1.263040753)
  .. (0.210000000,1.265511948)
  .. controls (0.213333333,1.267948067)
  and (0.216666667,1.270384187)
  .. (0.220000000,1.272820307)
  .. controls (0.223333333,1.275196186)
  and (0.226666667,1.277572066)
  .. (0.230000000,1.279947946)
  .. controls (0.233333333,1.282234071)
  and (0.236666667,1.284520195)
  .. (0.240000000,1.286806320)
  .. controls (0.243333333,1.288969371)
  and (0.246666667,1.291132421)
  .. (0.250000000,1.293295472)
  .. controls (0.253333333,1.295298980)
  and (0.256666667,1.297302487)
  .. (0.260000000,1.299305995)
  .. controls (0.263333333,1.301111096)
  and (0.266666667,1.302916197)
  .. (0.270000000,1.304721299)
  .. controls (0.273333333,1.306287571)
  and (0.276666667,1.307853843)
  .. (0.280000000,1.309420116)
  .. controls (0.283333333,1.310706480)
  and (0.286666667,1.311992844)
  .. (0.290000000,1.313279207)
  .. controls (0.293333333,1.314244872)
  and (0.296666667,1.315210537)
  .. (0.300000000,1.316176201)
  .. controls (0.303333333,1.316781637)
  and (0.306666667,1.317387073)
  .. (0.310000000,1.317992509)
  .. controls (0.313333333,1.318200423)
  and (0.316666667,1.318408337)
  .. (0.320000000,1.318616251)
  .. controls (0.323333333,1.318392545)
  and (0.326666667,1.318168840)
  .. (0.330000000,1.317945134)
  .. controls (0.333333333,1.317259826)
  and (0.336666667,1.316574517)
  .. (0.340000000,1.315889209)
  .. controls (0.343333333,1.314717289)
  and (0.346666667,1.313545368)
  .. (0.350000000,1.312373448)
  .. controls (0.353333333,1.310695662)
  and (0.356666667,1.309017877)
  .. (0.360000000,1.307340091)
  .. controls (0.363333333,1.305143626)
  and (0.366666667,1.302947160)
  .. (0.370000000,1.300750694)
  .. controls (0.373333333,1.298029742)
  and (0.376666667,1.295308790)
  .. (0.380000000,1.292587839)
  .. controls (0.383333333,1.289344046)
  and (0.386666667,1.286100253)
  .. (0.390000000,1.282856460)
  .. controls (0.393333333,1.279099225)
  and (0.396666667,1.275341990)
  .. (0.400000000,1.271584754)
  .. controls (0.403333333,1.267331385)
  and (0.406666667,1.263078015)
  .. (0.410000000,1.258824645)
  .. controls (0.413333333,1.254100359)
  and (0.416666667,1.249376073)
  .. (0.420000000,1.244651786)
  .. controls (0.423333333,1.239489555)
  and (0.426666667,1.234327324)
  .. (0.430000000,1.229165093)
  .. controls (0.433333333,1.223605334)
  and (0.436666667,1.218045574)
  .. (0.440000000,1.212485815)
  .. controls (0.443333333,1.206575922)
  and (0.446666667,1.200666029)
  .. (0.450000000,1.194756137)
  .. controls (0.453333333,1.188549878)
  and (0.456666667,1.182343620)
  .. (0.460000000,1.176137361)
  .. controls (0.463333333,1.169694132)
  and (0.466666667,1.163250902)
  .. (0.470000000,1.156807673)
  .. controls (0.473333333,1.150191628)
  and (0.476666667,1.143575584)
  .. (0.480000000,1.136959539)
  .. controls (0.483333333,1.130238624)
  and (0.486666667,1.123517709)
  .. (0.490000000,1.116796794)
  .. controls (0.493333333,1.110041678)
  and (0.496666667,1.103286562)
  .. (0.500000000,1.096531446)
  .. controls (0.503333333,1.089814392)
  and (0.506666667,1.083097338)
  .. (0.510000000,1.076380285)
  .. controls (0.513333333,1.069773970)
  and (0.516666667,1.063167655)
  .. (0.520000000,1.056561340)
  .. controls (0.523333333,1.050137647)
  and (0.526666667,1.043713954)
  .. (0.530000000,1.037290261)
  .. controls (0.533333333,1.031119070)
  and (0.536666667,1.024947879)
  .. (0.540000000,1.018776688)
  .. controls (0.543333333,1.012924688)
  and (0.546666667,1.007072687)
  .. (0.550000000,1.001220687)
  .. controls (0.553333333,0.995750231)
  and (0.556666667,0.990279774)
  .. (0.560000000,0.984809318)
  .. controls (0.563333333,0.979777341)
  and (0.566666667,0.974745365)
  .. (0.570000000,0.969713389)
  .. controls (0.573333333,0.965170422)
  and (0.576666667,0.960627455)
  .. (0.580000000,0.956084488)
  .. controls (0.583333333,0.952073767)
  and (0.586666667,0.948063045)
  .. (0.590000000,0.944052324)
  .. controls (0.593333333,0.940609031)
  and (0.596666667,0.937165738)
  .. (0.600000000,0.933722444)
  .. controls (0.603333333,0.930873088)
  and (0.606666667,0.928023731)
  .. (0.610000000,0.925174375)
  .. controls (0.613333333,0.922936321)
  and (0.616666667,0.920698267)
  .. (0.620000000,0.918460212)
  .. controls (0.623333333,0.916841377)
  and (0.626666667,0.915222541)
  .. (0.630000000,0.913603706)
  .. controls (0.633333333,0.912602419)
  and (0.636666667,0.911601133)
  .. (0.640000000,0.910599846)
  .. controls (0.643333333,0.910204886)
  and (0.646666667,0.909809926)
  .. (0.650000000,0.909414966)
  .. controls (0.653333333,0.909605764)
  and (0.656666667,0.909796563)
  .. (0.660000000,0.909987362)
  .. controls (0.663333333,0.910734381)
  and (0.666666667,0.911481400)
  .. (0.670000000,0.912228419)
  .. controls (0.673333333,0.913493686)
  and (0.676666667,0.914758953)
  .. (0.680000000,0.916024220)
  .. controls (0.683333333,0.917762019)
  and (0.686666667,0.919499818)
  .. (0.690000000,0.921237617)
  .. controls (0.693333333,0.923395315)
  and (0.696666667,0.925553013)
  .. (0.700000000,0.927710712)
  .. controls (0.703333333,0.930229713)
  and (0.706666667,0.932748713)
  .. (0.710000000,0.935267714)
  .. controls (0.713333333,0.938084516)
  and (0.716666667,0.940901317)
  .. (0.720000000,0.943718119)
  .. controls (0.723333333,0.946765457)
  and (0.726666667,0.949812795)
  .. (0.730000000,0.952860133)
  .. controls (0.733333333,0.956068190)
  and (0.736666667,0.959276247)
  .. (0.740000000,0.962484304)
  .. controls (0.743333333,0.965781957)
  and (0.746666667,0.969079610)
  .. (0.750000000,0.972377264)
  .. controls (0.753333333,0.975693355)
  and (0.756666667,0.979009445)
  .. (0.760000000,0.982325536)
  .. controls (0.763333333,0.985590129)
  and (0.766666667,0.988854722)
  .. (0.770000000,0.992119315)
  .. controls (0.773333333,0.995264929)
  and (0.776666667,0.998410543)
  .. (0.780000000,1.001556157)
  .. controls (0.783333333,1.004518944)
  and (0.786666667,1.007481730)
  .. (0.790000000,1.010444517)
  .. controls (0.793333333,1.013165362)
  and (0.796666667,1.015886207)
  .. (0.800000000,1.018607052)
  .. controls (0.803333333,1.021032582)
  and (0.806666667,1.023458113)
  .. (0.810000000,1.025883643)
  .. controls (0.813333333,1.027967120)
  and (0.816666667,1.030050596)
  .. (0.820000000,1.032134072)
  .. controls (0.823333333,1.033836150)
  and (0.826666667,1.035538228)
  .. (0.830000000,1.037240306)
  .. controls (0.833333333,1.038529654)
  and (0.836666667,1.039819002)
  .. (0.840000000,1.041108349)
  .. controls (0.843333333,1.041962110)
  and (0.846666667,1.042815870)
  .. (0.850000000,1.043669630)
  .. controls (0.853333333,1.044073722)
  and (0.856666667,1.044477813)
  .. (0.860000000,1.044881904)
  .. controls (0.863333333,1.044831152)
  and (0.866666667,1.044780400)
  .. (0.870000000,1.044729647)
  .. controls (0.873333333,1.044227750)
  and (0.876666667,1.043725852)
  .. (0.880000000,1.043223955)
  .. controls (0.883333333,1.042283281)
  and (0.886666667,1.041342607)
  .. (0.890000000,1.040401933)
  .. controls (0.893333333,1.039043160)
  and (0.896666667,1.037684388)
  .. (0.900000000,1.036325616)
  .. controls (0.903333333,1.034577219)
  and (0.906666667,1.032828821)
  .. (0.910000000,1.031080424)
  .. controls (0.913333333,1.028978017)
  and (0.916666667,1.026875609)
  .. (0.920000000,1.024773202)
  .. controls (0.923333333,1.022358759)
  and (0.926666667,1.019944317)
  .. (0.930000000,1.017529875)
  .. controls (0.933333333,1.014850845)
  and (0.936666667,1.012171815)
  .. (0.940000000,1.009492785)
  .. controls (0.943333333,1.006601104)
  and (0.946666667,1.003709424)
  .. (0.950000000,1.000817744)
  .. controls (0.953333333,0.997768791)
  and (0.956666667,0.994719838)
  .. (0.960000000,0.991670885)
  .. controls (0.963333333,0.988522377)
  and (0.966666667,0.985373869)
  .. (0.970000000,0.982225361)
  .. controls (0.973333333,0.979036230)
  and (0.976666667,0.975847100)
  .. (0.980000000,0.972657969)
  .. controls (0.983333333,0.969487232)
  and (0.986666667,0.966316494)
  .. (0.990000000,0.963145757)
  .. controls (0.993333333,0.960051403)
  and (0.996666667,0.956957049)
  .. (1.000000000,0.953862695)
  .. controls (1.003333333,0.950900617)
  and (1.006666667,0.947938539)
  .. (1.010000000,0.944976461)
  .. controls (1.013333333,0.942199449)
  and (1.016666667,0.939422436)
  .. (1.020000000,0.936645423)
  .. controls (1.023333333,0.934102232)
  and (1.026666667,0.931559042)
  .. (1.030000000,0.929015851)
  .. controls (1.033333333,0.926750380)
  and (1.036666667,0.924484910)
  .. (1.040000000,0.922219439)
  .. controls (1.043333333,0.920270014)
  and (1.046666667,0.918320589)
  .. (1.050000000,0.916371164)
  .. controls (1.053333333,0.914769965)
  and (1.056666667,0.913168766)
  .. (1.060000000,0.911567567)
  .. controls (1.063333333,0.910340278)
  and (1.066666667,0.909112993)
  .. (1.070000000,0.907885712)
  .. controls (1.073333333,0.907051564)
  and (1.076666667,0.906217432)
  .. (1.080000000,0.905383315)
  .. controls (1.083333333,0.904955558)
  and (1.086666667,0.904527831)
  .. (1.090000000,0.904100133)
  .. controls (1.093333333,0.904086671)
  and (1.096666667,0.904073214)
  .. (1.100000000,0.904059760)
  .. controls (1.103333333,0.904463526)
  and (1.106666667,0.904867132)
  .. (1.110000000,0.905270578)
  .. controls (1.113333333,0.906089411)
  and (1.116666667,0.906907618)
  .. (1.120000000,0.907725198)
  .. controls (1.123333333,0.908951262)
  and (1.126666667,0.910175723)
  .. (1.130000000,0.911398583)
  .. controls (1.133333333,0.913017546)
  and (1.136666667,0.914633213)
  .. (1.140000000,0.916245585)
  .. controls (1.143333333,0.918235880)
  and (1.146666667,0.920220295)
  .. (1.150000000,0.922198830)
  .. controls (1.153333333,0.924531231)
  and (1.156666667,0.926854175)
  .. (1.160000000,0.929167662)
  .. controls (1.163333333,0.931805356)
  and (1.166666667,0.934429012)
  .. (1.170000000,0.937038629)
  .. controls (1.173333333,0.939937837)
  and (1.176666667,0.942817510)
  .. (1.180000000,0.945677649)
  .. controls (1.183333333,0.948788784)
  and (1.186666667,0.951874151)
  .. (1.190000000,0.954933750)
  .. controls (1.193333333,0.958203024)
  and (1.196666667,0.961439804)
  .. (1.200000000,0.964644090)
  .. controls (1.203333333,0.968015424)
  and (1.206666667,0.971347338)
  .. (1.210000000,0.974639831)
  .. controls (1.213333333,0.978056926)
  and (1.216666667,0.981427778)
  .. (1.220000000,0.984752387)
  .. controls (1.223333333,0.988160785)
  and (1.226666667,0.991516512)
  .. (1.230000000,0.994819567)
  .. controls (1.233333333,0.998168561)
  and (1.236666667,1.001459102)
  .. (1.240000000,1.004691188)
  .. controls (1.243333333,1.007935462)
  and (1.246666667,1.011116341)
  .. (1.250000000,1.014233823)
  .. controls (1.253333333,1.017334725)
  and (1.256666667,1.020368255)
  .. (1.260000000,1.023334414)
  .. controls (1.263333333,1.026260816)
  and (1.266666667,1.029116891)
  .. (1.270000000,1.031902640)
  .. controls (1.273333333,1.034631372)
  and (1.276666667,1.037287825)
  .. (1.280000000,1.039871998)
  .. controls (1.283333333,1.042387863)
  and (1.286666667,1.044830416)
  .. (1.290000000,1.047199659)
  .. controls (1.293333333,1.049495065)
  and (1.296666667,1.051716919)
  .. (1.300000000,1.053865222)
  .. controls (1.303333333,1.055939508)
  and (1.306666667,1.057940624)
  .. (1.310000000,1.059868572)
  .. controls (1.313333333,1.061727077)
  and (1.316666667,1.063513230)
  .. (1.320000000,1.065227031)
  .. controls (1.323333333,1.066880001)
  and (1.326666667,1.068461675)
  .. (1.330000000,1.069972052)
  .. controls (1.333333333,1.071433447)
  and (1.336666667,1.072824651)
  .. (1.340000000,1.074145665)
  .. controls (1.343333333,1.075431943)
  and (1.346666667,1.076649018)
  .. (1.350000000,1.077796891)
  .. controls (1.353333333,1.078925831)
  and (1.356666667,1.079986299)
  .. (1.360000000,1.080978293)
  .. controls (1.363333333,1.081967915)
  and (1.366666667,1.082889427)
  .. (1.370000000,1.083742830)
  .. controls (1.373333333,1.084610437)
  and (1.376666667,1.085409864)
  .. (1.380000000,1.086141112)
  .. controls (1.383333333,1.086902503)
  and (1.386666667,1.087595181)
  .. (1.390000000,1.088219147)
  .. controls (1.393333333,1.088888008)
  and (1.396666667,1.089487169)
  .. (1.400000000,1.090016632)
  .. controls (1.403333333,1.090604113)
  and (1.406666667,1.091120500)
  .. (1.410000000,1.091565793)
  .. controls (1.413333333,1.092080281)
  and (1.416666667,1.092521946)
  .. (1.420000000,1.092890786)
  .. controls (1.423333333,1.093337847)
  and (1.426666667,1.093710123)
  .. (1.430000000,1.094007612)
  .. controls (1.433333333,1.094390105)
  and (1.436666667,1.094695737)
  .. (1.440000000,1.094924511)
  .. controls (1.443333333,1.095242838)
  and (1.446666667,1.095482254)
  .. (1.450000000,1.095642759)
  .. controls (1.453333333,1.095895247)
  and (1.456666667,1.096066934)
  .. (1.460000000,1.096157819)
  .. controls (1.463333333,1.096341181)
  and (1.466666667,1.096442154)
  .. (1.470000000,1.096460736)
  .. controls (1.473333333,1.096570607)
  and (1.476666667,1.096596933)
  .. (1.480000000,1.096539716)
  .. controls (1.483333333,1.096571216)
  and (1.486666667,1.096518575)
  .. (1.490000000,1.096381793)
  .. controls (1.493333333,1.096330107)
  and (1.496666667,1.096194343)
  .. (1.500000000,1.095974503)
  .. controls (1.503333333,1.095835432)
  and (1.506666667,1.095613093)
  .. (1.510000000,1.095307485)
  .. controls (1.513333333,1.095077959)
  and (1.516666667,1.094766777)
  .. (1.520000000,1.094373937)
  .. controls (1.523333333,1.094052458)
  and (1.526666667,1.093651767)
  .. (1.530000000,1.093171865)
  .. controls (1.533333333,1.092758857)
  and (1.536666667,1.092269921)
  .. (1.540000000,1.091705057)
  .. controls (1.543333333,1.091203125)
  and (1.546666667,1.090629358)
  .. (1.550000000,1.089983757)
  .. controls (1.553333333,1.089397822)
  and (1.556666667,1.088744901)
  .. (1.560000000,1.088024992)
  .. controls (1.563333333,1.087362315)
  and (1.566666667,1.086638169)
  .. (1.570000000,1.085852552)
  .. controls (1.573333333,1.085122626)
  and (1.576666667,1.084337308)
  .. (1.580000000,1.083496598)
  .. controls (1.583333333,1.082710933)
  and (1.586666667,1.081876382)
  .. (1.590000000,1.080992945)
  .. controls (1.593333333,1.080164734)
  and (1.596666667,1.079294424)
  .. (1.600000000,1.078382016)
  .. controls (1.603333333,1.077525715)
  and (1.606666667,1.076634223)
  .. (1.610000000,1.075707539)
  .. controls (1.613333333,1.074838371)
  and (1.616666667,1.073940871)
  .. (1.620000000,1.073015038)
  .. controls (1.623333333,1.072148437)
  and (1.626666667,1.071260150)
  .. (1.630000000,1.070350176)
  .. controls (1.633333333,1.069501217)
  and (1.636666667,1.068636847)
  .. (1.640000000,1.067757065)
  .. controls (1.643333333,1.066939894)
  and (1.646666667,1.066113071)
  .. (1.650000000,1.065276595)
  .. controls (1.653333333,1.064503897)
  and (1.656666667,1.063726666)
  .. (1.660000000,1.062944903)
  .. controls (1.663333333,1.062227432)
  and (1.666666667,1.061509816)
  .. (1.670000000,1.060792053)
  .. controls (1.673333333,1.060138263)
  and (1.676666667,1.059487917)
  .. (1.680000000,1.058841015)
  .. controls (1.683333333,1.058256803)
  and (1.686666667,1.057678799)
  .. (1.690000000,1.057107004)
  .. controls (1.693333333,1.056595590)
  and (1.696666667,1.056092332)
  .. (1.700000000,1.055597229)
  .. controls (1.703333333,1.055159183)
  and (1.706666667,1.054730466)
  .. (1.710000000,1.054311076)
  .. controls (1.713333333,1.053944483)
  and (1.716666667,1.053587691)
  .. (1.720000000,1.053240702)
  .. controls (1.723333333,1.052941460)
  and (1.726666667,1.052651897)
  .. (1.730000000,1.052372013)
  .. controls (1.733333333,1.052134247)
  and (1.736666667,1.051905560)
  .. (1.740000000,1.051685955)
  .. controls (1.743333333,1.051502499)
  and (1.746666667,1.051327181)
  .. (1.750000000,1.051160001)
  .. controls (1.753333333,1.051022934)
  and (1.756666667,1.050892851)
  .. (1.760000000,1.050769751)
  .. controls (1.763333333,1.050670919)
  and (1.766666667,1.050577828)
  .. (1.770000000,1.050490477)
  .. controls (1.773333333,1.050421979)
  and (1.776666667,1.050357996)
  .. (1.780000000,1.050298528)
  .. controls (1.783333333,1.050253114)
  and (1.786666667,1.050211090)
  .. (1.790000000,1.050172455)
  .. controls (1.793333333,1.050143822)
  and (1.796666667,1.050117609)
  .. (1.800000000,1.050093816)
  .. controls (1.803333333,1.050076770)
  and (1.806666667,1.050061360)
  .. (1.810000000,1.050047587)
  .. controls (1.813333333,1.050038091)
  and (1.816666667,1.050029635)
  .. (1.820000000,1.050022220)
  .. controls (1.823333333,1.050017326)
  and (1.826666667,1.050013049)
  .. (1.830000000,1.050009389)
  .. controls (1.833333333,1.050007090)
  and (1.836666667,1.050005128)
  .. (1.840000000,1.050003505)
  .. controls (1.843333333,1.050002539)
  and (1.846666667,1.050001742)
  .. (1.850000000,1.050001114)
  .. controls (1.853333333,1.050000761)
  and (1.856666667,1.050000484)
  .. (1.860000000,1.050000283)
  .. controls (1.863333333,1.050000175)
  and (1.866666667,1.050000097)
  .. (1.870000000,1.050000049)
  .. controls (1.873333333,1.050000024)
  and (1.876666667,1.050000008)
  .. (1.880000000,1.050000002)
  .. controls (1.883333333,1.049999998)
  and (1.886666667,1.049999997)
  .. (1.890000000,1.049999998)
  .. controls (1.893333333,1.049999998)
  and (1.896666667,1.049999998)
  .. (1.900000000,1.049999999)
  .. controls (1.903333333,1.049999999)
  and (1.906666667,1.050000000)
  .. (1.910000000,1.050000000)
  .. controls (1.913333333,1.050000000)
  and (1.916666667,1.050000000)
  .. (1.920000000,1.050000000)
  .. controls (1.923333333,1.050000000)
  and (1.926666667,1.050000000)
  .. (1.930000000,1.050000000)
  .. controls (1.933333333,1.050000000)
  and (1.936666667,1.050000000)
  .. (1.940000000,1.050000000)
  .. controls (1.943333333,1.050000000)
  and (1.946666667,1.050000000)
  .. (1.950000000,1.050000000)
  .. controls (1.953333333,1.050000000)
  and (1.956666667,1.050000000)
  .. (1.960000000,1.050000000)
  .. controls (1.963333333,1.050000000)
  and (1.966666667,1.050000000)
  .. (1.970000000,1.050000000)
  .. controls (1.973333333,1.050000000)
  and (1.976666667,1.050000000)
  .. (1.980000000,1.050000000)
  .. controls (1.983333333,1.050000000)
  and (1.986666667,1.050000000)
  .. (1.990000000,1.050000000)
  .. controls (1.993333333,1.050000000)
  and (1.996666667,1.050000000)
  .. (2.000000000,1.050000000)
  .. controls (2.003333333,1.050000000)
  and (2.006666667,1.050000000)
  .. (2.010000000,1.050000000)
  .. controls (2.013333333,1.050000000)
  and (2.016666667,1.050000000)
  .. (2.020000000,1.050000000)
  .. controls (2.023333333,1.050000000)
  and (2.026666667,1.050000000)
  .. (2.030000000,1.050000000)
  .. controls (2.033333333,1.050000000)
  and (2.036666667,1.050000000)
  .. (2.040000000,1.050000000)
  .. controls (2.043333333,1.050000000)
  and (2.046666667,1.050000000)
  .. (2.050000000,1.050000000)
  .. controls (2.053333333,1.050000000)
  and (2.056666667,1.050000000)
  .. (2.060000000,1.050000000)
  .. controls (2.063333333,1.050000000)
  and (2.066666667,1.050000000)
  .. (2.070000000,1.050000000)
  .. controls (2.073333333,1.050000000)
  and (2.076666667,1.050000000)
  .. (2.080000000,1.050000000)
  .. controls (2.083333333,1.050000000)
  and (2.086666667,1.050000000)
  .. (2.090000000,1.050000000)
  .. controls (2.093333333,1.050000000)
  and (2.096666667,1.050000000)
  .. (2.100000000,1.050000000)
  .. controls (2.103333333,1.050000000)
  and (2.106666667,1.050000000)
  .. (2.110000000,1.050000000)
  .. controls (2.113333333,1.050000000)
  and (2.116666667,1.050000000)
  .. (2.120000000,1.050000000)
  .. controls (2.123333333,1.050000000)
  and (2.126666667,1.050000000)
  .. (2.130000000,1.050000000)
  .. controls (2.133333333,1.050000000)
  and (2.136666667,1.050000000)
  .. (2.140000000,1.050000000)
  .. controls (2.143333333,1.050000000)
  and (2.146666667,1.050000000)
  .. (2.150000000,1.050000000)
  .. controls (2.153333333,1.050000000)
  and (2.156666667,1.050000000)
  .. (2.160000000,1.050000000)
  .. controls (2.163333333,1.050000000)
  and (2.166666667,1.050000000)
  .. (2.170000000,1.050000000)
  .. controls (2.173333333,1.050000000)
  and (2.176666667,1.050000000)
  .. (2.180000000,1.050000000)
  .. controls (2.183333333,1.050000000)
  and (2.186666667,1.050000000)
  .. (2.190000000,1.050000000)
  .. controls (2.193333333,1.050000000)
  and (2.196666667,1.050000000)
  .. (2.200000000,1.050000000)
  .. controls (2.203333333,1.050000000)
  and (2.206666667,1.050000000)
  .. (2.210000000,1.050000000)
  .. controls (2.213333333,1.050000000)
  and (2.216666667,1.050000000)
  .. (2.220000000,1.050000000)
  .. controls (2.223333333,1.050000000)
  and (2.226666667,1.050000000)
  .. (2.230000000,1.050000000)
  .. controls (2.233333333,1.050000000)
  and (2.236666667,1.050000000)
  .. (2.240000000,1.050000000)
  .. controls (2.243333333,1.050000000)
  and (2.246666667,1.050000000)
  .. (2.250000000,1.050000000)
  .. controls (2.253333333,1.050000000)
  and (2.256666667,1.050000000)
  .. (2.260000000,1.050000000)
  .. controls (2.263333333,1.050000000)
  and (2.266666667,1.050000000)
  .. (2.270000000,1.050000000)
  .. controls (2.273333333,1.050000000)
  and (2.276666667,1.050000000)
  .. (2.280000000,1.050000000)
  .. controls (2.283333333,1.050000000)
  and (2.286666667,1.050000000)
  .. (2.290000000,1.050000000)
  .. controls (2.293333333,1.050000000)
  and (2.296666667,1.050000000)
  .. (2.300000000,1.050000000)
  .. controls (2.303333333,1.050000000)
  and (2.306666667,1.050000000)
  .. (2.310000000,1.050000000)
  .. controls (2.313333333,1.050000000)
  and (2.316666667,1.050000000)
  .. (2.320000000,1.050000000)
  .. controls (2.323333333,1.050000000)
  and (2.326666667,1.050000000)
  .. (2.330000000,1.050000000)
  .. controls (2.333333333,1.050000000)
  and (2.336666667,1.050000000)
  .. (2.340000000,1.050000000)
  .. controls (2.343333333,1.050000000)
  and (2.346666667,1.050000000)
  .. (2.350000000,1.050000000)
  .. controls (2.353333333,1.050000000)
  and (2.356666667,1.050000000)
  .. (2.360000000,1.050000000)
  .. controls (2.363333333,1.050000000)
  and (2.366666667,1.050000000)
  .. (2.370000000,1.050000000)
  .. controls (2.373333333,1.050000000)
  and (2.376666667,1.050000000)
  .. (2.380000000,1.050000000)
  .. controls (2.383333333,1.050000000)
  and (2.386666667,1.050000000)
  .. (2.390000000,1.050000000)
  .. controls (2.393333333,1.050000000)
  and (2.396666667,1.050000000)
  .. (2.400000000,1.050000000)
  .. controls (2.403333333,1.050000000)
  and (2.406666667,1.050000000)
  .. (2.410000000,1.050000000)
  .. controls (2.413333333,1.050000000)
  and (2.416666667,1.050000000)
  .. (2.420000000,1.050000000)
  .. controls (2.423333333,1.050000000)
  and (2.426666667,1.050000000)
  .. (2.430000000,1.050000000)
  .. controls (2.433333333,1.050000000)
  and (2.436666667,1.050000000)
  .. (2.440000000,1.050000000)
  .. controls (2.443333333,1.050000000)
  and (2.446666667,1.050000000)
  .. (2.450000000,1.050000000)
  .. controls (2.453333333,1.050000000)
  and (2.456666667,1.050000000)
  .. (2.460000000,1.050000000)
  .. controls (2.463333333,1.050000000)
  and (2.466666667,1.050000000)
  .. (2.470000000,1.050000000)
  .. controls (2.473333333,1.050000000)
  and (2.476666667,1.050000000)
  .. (2.480000000,1.050000000)
  .. controls (2.483333333,1.050000000)
  and (2.486666667,1.050000000)
  .. (2.490000000,1.050000000)
  .. controls (2.493333333,1.050000000)
  and (2.496666667,1.050000000)
  .. (2.500000000,1.050000000)
  .. controls (2.503333333,1.050000000)
  and (2.506666667,1.050000000)
  .. (2.510000000,1.050000000)
  .. controls (2.513333333,1.050000000)
  and (2.516666667,1.050000000)
  .. (2.520000000,1.050000000)
  .. controls (2.523333333,1.050000000)
  and (2.526666667,1.050000000)
  .. (2.530000000,1.050000000)
  .. controls (2.533333333,1.050000000)
  and (2.536666667,1.050000000)
  .. (2.540000000,1.050000000)
  .. controls (2.543333333,1.050000000)
  and (2.546666667,1.050000000)
  .. (2.550000000,1.050000000)
  .. controls (2.553333333,1.050000000)
  and (2.556666667,1.050000000)
  .. (2.560000000,1.050000000)
  .. controls (2.563333333,1.050000000)
  and (2.566666667,1.050000000)
  .. (2.570000000,1.050000000)
  .. controls (2.573333333,1.050000000)
  and (2.576666667,1.050000000)
  .. (2.580000000,1.050000000)
  .. controls (2.583333333,1.050000000)
  and (2.586666667,1.050000000)
  .. (2.590000000,1.050000000)
  .. controls (2.593333333,1.050000000)
  and (2.596666667,1.050000000)
  .. (2.600000000,1.050000000)
  .. controls (2.603333333,1.050000000)
  and (2.606666667,1.050000000)
  .. (2.610000000,1.050000000)
  .. controls (2.613333333,1.050000000)
  and (2.616666667,1.050000000)
  .. (2.620000000,1.050000000)
  .. controls (2.623333333,1.050000000)
  and (2.626666667,1.050000000)
  .. (2.630000000,1.050000000)
  .. controls (2.633333333,1.050000000)
  and (2.636666667,1.050000000)
  .. (2.640000000,1.050000000)
  .. controls (2.643333333,1.050000000)
  and (2.646666667,1.050000000)
  .. (2.650000000,1.050000000);
  \draw[ham/arrow] (-2.7,.35)--(2.65,.35) node[right] {$x$};
  \foreach \xx/\tick in {-2/-2R,-1/-R,1/R,2/2R} {
  \draw[hamblue] (\xx,.30)--(\xx,.40);
  \node[anchor=north,inner sep=1pt] at (\xx,.22) {$\tick$};
  }
\end{tikzpicture}

  \caption[A physical mask as a window on the potential]{A physical mask,
    equal to one on a ball and zero outside a larger ball, and its product
    with the potential. The mask multiplies the potential in space. The
    covariance of the product has a mask factor at each spatial argument. The
    original noise and the covariance norm used for control costs remain
    fixed. A schematic in one space dimension: the unsmoothed potential is a
    distribution, not a curve.}
  \label{Fig:mask-window}
\end{figure}


We establish the covariance estimates needed for the operator construction.
Throughout this section, Hypothesis~\ref{H:homog} holds, $a = \alpha_0$, and
$\Gamma \ne 0$. For zero spatial covariance, the spatial spectral measures and
noise variances vanish; the purely temporal estimates remain valid independently
of $\Gamma$. We use the unitary Fourier transform
\[
  \widehat{f}(\xi) = c_d\int_{\R^d}e^{-i\xi\cdot x}f(x)\ud x,
  \qquad c_d = (2\pi)^{-d/2}.
\]
The spectral convention in Hypothesis~\ref{H:homog} means that, for every
complex Schwartz test $\varphi$,
\begin{equation}\label{E:masked-spectral-convention}
  \int_{\R^d}\varphi(x)\Gamma(\ud x)
  = \int_{\R^d}\left(\int_{\R^d}e^{i\xi\cdot x}
  \varphi(x)\ud x\right)\mu(\ud\xi).
\end{equation}
This is an identity of distribution pairings. We will not interchange two
infinite measures in the associated oscillatory expression.

A positive tempered Radon measure has polynomial mass growth. Indeed, choose a
nonnegative compact smooth test that is at least one on the unit ball, and apply
a defining Schwartz seminorm bound for the distribution to its dilations. Its
mass on a ball of radius $R \ge 1$ is bounded by a power of $R$. Summation over
dyadic shells then gives integrability of $(1+|x|)^{-N}$ for some $N$, and hence
absolute integrability of every Schwartz test. These facts apply to both
$\Gamma$ and $\mu$. Smooth cutoffs converge in the Schwartz topology, so these
integrals agree with the tempered pairings.

\subsection{Translated spectral growth and physical masks}

\begin{proposition}[Translated spectral mass]\label{P:masked-spectral-growth}
  Put
  \[
    M_\Gamma = \int_{\R^d}e^{-|x|^2/4}\Gamma(\ud x),
    \qquad C_\Gamma = \frac{eM_\Gamma}{(2\sqrt{\pi})^d}.
  \]
  Then $0 < M_\Gamma < \infty$, and for $y \in \R^d$ and $r > 0$,
  \begin{align}
    \int_{\R^d}e^{-|\xi-y|^2/r^2}\mu(\ud\xi)
    &= \left(\frac{r}{2\sqrt{\pi}}\right)^d
    \int_{\R^d}e^{-r^2|x|^2/4}e^{-iy\cdot x}\Gamma(\ud x),
    \label{E:masked-gaussian-transform}\\
    0 \le \int_{\R^d}e^{-|\xi-y|^2/r^2}\mu(\ud\xi)
    &\le \frac{r^\alpha M_\Gamma}{(2\sqrt{\pi})^d},
    \qquad \mu(B(y,r)) \le C_\Gamma r^\alpha.
    \label{E:masked-translated-ball}
  \end{align}
  The ball may be open or closed.
\end{proposition}

\begin{proof}
  Temperedness gives finiteness of $M_\Gamma$. Its positivity follows from
  $\Gamma \ne 0$. To check the Fourier constant, set $J(b) = \int_\R
  e^{-Ax^2}e^{ibx}\ud x$ for $A > 0$. Squaring the real Gaussian integral and
  using polar coordinates gives $J(0) = \sqrt{\pi/A}$. Differentiation and
  integration by parts give $J'(b) = -bJ(b)/(2A)$. Therefore, $J(b) =
  \sqrt{\pi/A}e^{-b^2/(4A)}$. Taking products over the coordinates gives
  \[
    \int_{\R^d}e^{-r^2|x|^2/4}e^{i(\xi-y)\cdot x}\ud x
    = \left(\frac{2\sqrt{\pi}}r\right)^d e^{-|\xi-y|^2/r^2}.
  \]
  Apply~\eqref{E:masked-spectral-convention} to the single Schwartz test
  $e^{-r^2|x|^2/4}e^{-iy\cdot x}$ to obtain \eqref{E:masked-gaussian-transform}.
  Its left side is nonnegative. Taking the absolute value on the right removes
  the complex phase because $\Gamma$ is nonnegative. The pushforward of $\Gamma$
  by $x \mapsto rx$ assigns mass $r^{\alpha-d}\Gamma(A)$ to a Borel set $A$.
  Thus,
  \[
    \int e^{-r^2|x|^2/4}\Gamma(\ud x) = r^{\alpha-d}M_\Gamma.
  \]
  This proves the Gaussian bound. On either ball $B(y,r)$ the integrand
  $e^{-|\xi-y|^2/r^2}$ is at least $e^{-1}$, proving the mass bound. This proves
  Proposition~\ref{P:masked-spectral-growth}.
\end{proof}

Choose and fix a real nonzero Schwartz function $\zeta$. For $R > 0$ put
$\zeta_R(x) = \zeta(x/R)$, $k_R = \widehat{\zeta_R}$ and $L =
\|\widehat{\zeta}\|_1$. Then $k_R(u) = R^d\widehat{\zeta}(Ru)$ and $\|k_R\|_1 =
L > 0$. Define
\begin{equation}\label{E:masked-auxiliary-measure}
  \nu_R(\ud\eta) = \rho_R(\eta)\ud\eta,
  \qquad \rho_R(\eta) = c_d^2L\int |k_R(\xi-\eta)|\mu(\ud\xi).
\end{equation}

\begin{proposition}[Auxiliary mask norm]\label{P:masked-auxiliary-norm}
  The measure $\nu_R$ is even, tempered and Radon, with
  \begin{equation}\label{E:masked-auxiliary-growth}
    \nu_R(B(y,r)) \le c_d^2L^2C_\Gamma r^\alpha
  \end{equation}
  for all $y \in \R^d$ and $r > 0$, uniformly in $R > 0$. If $H$ is a separable
  complex Hilbert space and $F$ is an $H$-valued Schwartz function, then
  \begin{equation}\label{E:masked-hilbert-fourier}
    \int \|\widehat{\zeta_RF}(\xi)\|_H^2\mu(\ud\xi)
    \le \int \|\widehat{F}(\eta)\|_H^2\nu_R(\ud\eta).
  \end{equation}
  For a real $H$, the same assertion holds after complexification.
\end{proposition}

\begin{proof}
  Schwartz decay of $k_R$ and polynomial growth of $\mu$ show that
  $\rho_R(\eta)$ is finite at every $\eta$. On compact sets of $\eta$, the
  integrands have a common integrable bound. Dominated convergence therefore
  makes $\rho_R$ continuous, and $\nu_R$ is Radon. Since $\zeta$ is real,
  $k_R(-u) = \overline{k_R(u)}$. Evenness of $\mu$ gives evenness of $\rho_R$.
  Tonelli's theorem gives
  \[
    \nu_R(B(y,r)) = c_d^2L\int |k_R(u)|\mu(B(y+u,r))\ud u.
  \]
  Proposition~\ref{P:masked-spectral-growth} and $\|k_R\|_1 = L$ give
  \eqref{E:masked-auxiliary-growth}. Polynomial growth on centered balls proves
  temperedness by the same dyadic argument as above.

  Bochner Fourier integration gives
  \[
    \widehat{\zeta_RF}(\xi)
    = c_d\int k_R(\xi-\eta)\widehat{F}(\eta)\ud\eta.
  \]
  For scalar Schwartz functions this follows by Fourier inversion and absolute
  integrability. It follows next for finite-rank Hilbert-valued functions. For
  general $F$, take projections onto the successive spans of a countable
  orthonormal basis of its separable range closure. They converge pointwise in
  norm. The Schwartz bounds on $F$ dominate the Fourier integrals, proving the
  identity by dominated convergence. In a finite-dimensional space the full
  identity projection suffices; the zero space is immediate. The triangle
  inequality and weighted Cauchy--Schwarz yield
  \[
    \|\widehat{\zeta_RF}(\xi)\|_H^2
    \le c_d^2L\int |k_R(\xi-\eta)|\|\widehat{F}(\eta)\|_H^2\ud\eta.
  \]
  Integrating against $\mu$ and applying Tonelli's theorem proves
  \eqref{E:masked-hilbert-fourier}. Its right side is finite by Schwartz decay
  and temperedness of $\nu_R$. This proves
  Proposition~\ref{P:masked-auxiliary-norm}.
\end{proof}

\subsection{Temporal completion and the original noise}

\begin{lemma}[Temporal covariance completion]\label{L:masked-time-completion}
  Let $T > 0$. The form
  \[
    \langle f,g\rangle_a
    = \int_0^T\int_0^T f(t)g(s)|t-s|^{-a}\ud t\ud s
  \]
  is positive semidefinite on real compact smooth tests. Its null-quotient
  completion $H_a$ is separable. The inclusion of tests extends continuously
  from $L^2(0,T)$, with
  \begin{equation}\label{E:masked-time-schur}
    \|f\|_a^2 \le \frac{2T^{1-a}}{1-a}\|f\|_2^2.
  \end{equation}
  In its complexification the norm is given by the same double integral with
  $f(t)\overline{f(s)}$. In particular, bounded measurable amplitudes on a
  subinterval, extended by zero, belong to this completion.
\end{lemma}

\begin{proof}
  Write $\Gamma_E$ for Euler's Gamma function. Away from the null diagonal,
  \[
    |t-s|^{-a} = \frac1{\Gamma_E(a)}\int_0^\infty
    \lambda^{a-1}e^{-\lambda|t-s|}\ud\lambda.
  \]
  This follows by substitution in the defining Gamma integral. Moreover,
  \[
    e^{-\lambda|t-s|}
    = \int_{-\infty}^{\min(t,s)}2\lambda
    e^{-\lambda(t-u)}e^{-\lambda(s-u)}\ud u.
  \]
  For each $\lambda > 0$ its quadratic form is an integral of a square. All
  exchanges for compact smooth tests are absolutely integrable because the
  kernel row integral is at most $2T^{1-a}/(1-a)$. The same bound and $2|xy| \le
  |x|^2+|y|^2$ prove~\eqref{E:masked-time-schur}. A countable $L^2$-dense family
  of compact smooth tests is consequently dense in $H_a$, proving separability.
  The complex quadratic form is the sum of those of the real and imaginary
  parts, since symmetry cancels the imaginary cross term. Thus, the same bound
  holds after complexification. Smooth density in $L^2$, together with the Schur
  bound on the absolute bilinear form, extends the inclusion and the integral
  identity to $L^2$. This includes the stated interval amplitudes. This proves
  Lemma~\ref{L:masked-time-completion}.
\end{proof}

Let $\mathcal{D}$ consist of finite sums $F(t,x) = \sum_j\tau_j(t)f_j(x)$ with
real $\tau_j \in C_c^\infty(0,T)$ and real Schwartz $f_j$. These give
$H_a$-valued Schwartz functions of $x$. Let $W$ denote the original real
isonormal noise in~\eqref{E:cov}.

\begin{proposition}[Completion on the same noise]\label{P:masked-noise-completion}
  The space $\mathcal{D}$ belongs to the original test-noise completion and
  \begin{equation}\label{E:masked-actual-noise-norm}
    \E W(F)^2 = (2\pi)^d\int
    \|\widehat{F}(\xi)\|_{H_a^{\mathbb{C}}}^2\mu(\ud\xi).
  \end{equation}
  Define the auxiliary seminorm
  \[
    N_R(F)^2 = (2\pi)^d\int
    \|\widehat{F}(\eta)\|_{H_a^{\mathbb{C}}}^2\nu_R(\ud\eta).
  \]
  Then $F \mapsto W(\zeta_RF)$ extends as a contraction from the null-quotient
  completion for $N_R$ into the original $L^2(\Omega)$. The extension is
  centered and jointly Gaussian. For finitely many real $F_1,\ldots,F_m \in
  \mathcal{D}$, their masked covariance matrix is bounded in positive
  semidefinite order by the matrix of auxiliary inner products. This gives
  covariance comparison for every finite linear Gaussian image and every convex
  norm power $p \ge 1$.
\end{proposition}

\begin{proof}
  For real spatial Schwartz $f,g$, put $h(z) = \int f(y+z)g(y)\ud y$. Fourier
  inversion and Parseval give
  \[
    h(z) = \int e^{i\xi\cdot z}\widehat{f}(\xi)
    \overline{\widehat{g}(\xi)}\ud\xi.
  \]
  Therefore,
  \[
    \int e^{i\xi\cdot z}h(z)\ud z
    = (2\pi)^d\widehat{f}(-\xi)\overline{\widehat{g}(-\xi)}.
  \]
  Apply~\eqref{E:masked-spectral-convention} to the single test $h$ and reflect
  $\xi$ using evenness of $\mu$. The spatial covariance is thus
  $(2\pi)^d\int\widehat{f}(\xi)\overline{\widehat{g}(\xi)}\mu(\ud\xi)$. Summing
  over the finitely many temporal factors gives
  \eqref{E:masked-actual-noise-norm}. Schwartz decay, temperedness and temporal
  integrability justify these pairings. If $W$ is initially given only on
  compact spatial tests, smooth expanding cutoffs converge in the Schwartz
  topology. The spectral form is continuous in that topology, admitting
  $\mathcal{D}$ in the same completion.

  Proposition~\ref{P:masked-auxiliary-norm}, applied with $H =
  H_a^{\mathbb{C}}$, gives $\E W(\zeta_RF)^2 \le N_R(F)^2$. It annihilates
  auxiliary-null tests. An auxiliary Cauchy sequence thus gives a Cauchy
  sequence of noise evaluations in $L^2(\Omega)$, proving the contraction
  extension and its uniqueness. Every finite linear combination of the limits is
  an $L^2$ limit of centered Gaussian variables. Their variances converge, so
  their characteristic functions prove that the limits remain jointly centered
  Gaussian.

  For the finite family let $A$ be the covariance matrix and $B$ the auxiliary
  Gram matrix. Applying the norm inequality to $\sum_i b_iF_i$ gives
  $b^{\mathsf{T}}Ab \le b^{\mathsf{T}}Bb$ for every real vector $b$.
  Diagonalizing $B-A$ constructs a centered Gaussian vector $Z$ of that
  covariance, independent of the evaluation vector $X$, on a product space. Then
  $X+Z$ has covariance $B$. Conditional Jensen gives $\E\|X\|^p \le \E\|X+Z\|^p$
  for every norm and $p \ge 1$. Finite Gaussian moments ensure finiteness of
  both sides. This proves Proposition~\ref{P:masked-noise-completion}.
\end{proof}

The auxiliary completion depends on $R$ and $\zeta$. Its construction does not
define multiplication by $\zeta_R$ on the original noise quotient. Nor does the
finite linear Gaussian comparison assert a covariance comparison for nonlinear
wave solutions. The auxiliary and original test norms need not dominate one
another. No spatial representative for every auxiliary class, nonnegative
physical inverse covariance for $\nu_R$, or uniform density bound in $R$ is
asserted.

\begin{proposition}[A physical mask need not preserve noise-null tests]
  \label{P:masked-quotient-obstruction}
  In the admissible case $d = 2$, $\alpha = 1$, $0 < a < 1$ and $\Gamma =
  \delta_0(\ud x_1)\ud x_2$, there is a real compact smooth mask $\zeta$ such
  that no finite $C$ satisfies $\E W(\zeta F)^2 \le C\E W(F)^2$ for every
  compact smooth test $F$.
\end{proposition}

\begin{proof}
  Fourier inversion on Schwartz tests gives $\mu =
  (2\pi)^{-1}\ud\xi_1\delta_0(\ud\xi_2)$ in the stated spectral convention. The
  measure $\Gamma$ is nonnegative, even, positive definite and tempered, with
  $\Gamma(cA) = c\Gamma(A)$. It therefore satisfies Hypothesis~\ref{H:homog}
  with $\alpha = 1$. This is not spatial white noise in dimension two. For real
  compact smooth $g$, its spatial covariance is, by ordinary Fubini,
  \[
    V(g) = \int_\R\left(\int_\R g(x_1,x_2)\ud x_2\right)^2\ud x_1.
  \]
  Choose real even $\chi \in C_c^\infty(\R)$ equal to one on $[-1,1]$ and zero
  outside $[-2,2]$. For example, on $1 < r = |x| < 2$ use
  \[
    \chi(x) = \frac{e^{-1/(2-r)}}{e^{-1/(2-r)}+e^{-1/(r-1)}}.
  \]
  The endpoint exponentials vanish to every order, so the plateau pastes
  smoothly. Let $I_\chi = \int(\chi')^2 > 0$ and $J_\chi = \int\chi^4 > 0$.
  Positivity of $I_\chi$ follows because $\chi$ is not constant, and $J_\chi \ge
  2$. Put $g(x) = \chi(x_1)\chi''(x_2)$ and $\zeta(x) = \chi(x_1)\chi(x_2)$.
  Their unmasked marginal vanishes since $\int\chi'' = 0$, whereas integration
  by parts gives the masked marginal $-I_\chi\chi(x_1)^2$. Thus, $V(g) = 0$ and
  $V(\zeta g) = J_\chi I_\chi^2 > 0$. For a nonzero nonnegative $\tau \in
  C_c^\infty(0,T)$, set $F(t,x) = \tau(t)g(x)$. The temporal factor
  $\|\tau\|_a^2$ is finite and strictly positive because its nonnegative
  integrand is positive on a set of positive two-dimensional measure.
  Consequently, the full unmasked variance is zero and the masked variance is
  positive. This proves Proposition~\ref{P:masked-quotient-obstruction}.
\end{proof}

No-go~\ref{P:nogo-22} records the quotient interpretation of this witness:
the mask separates two representatives of the same noise-null class.
This is why we mask tests before constructing the auxiliary completion.


\subsection{Temporal modulation and continuous wave factors}

Write $\langle x\rangle = \sqrt{1+x^2}$. Let $I = [s,s+h] \subset [0,T]$, with
$0 < h \le 1$. Norms below are those of the complexified temporal completion in
Lemma~\ref{L:masked-time-completion}.

\begin{lemma}[Modulated amplitudes]\label{L:masked-temporal-modulation}
  Let $A \in C^1(I;\mathbb{C})$, $M = \sup_I|A|$ and $D = \sup_I|A'|$. For every
  real $\omega$,
  \begin{equation}\label{E:masked-temporal-modulation}
    \|1_IA(t)e^{i\omega t}\|_a^2
    \le C_0(a)h^{2-a}(M+hD)^2\langle h\omega\rangle^{a-1},
    \qquad C_0(a) = \frac{16}{1-a}.
  \end{equation}
  The bound is uniform in $T,s,h,A$ and $\omega$.
\end{lemma}

\begin{proof}
  Put $c(v) = \int_s^{s+h-v}A(u+v)\overline{A(u)}\ud u$ for $0 \le v \le h$.
  Differentiation at the moving endpoint gives
  \[
    c'(v) = -A(s+h)\overline{A(s+h-v)}
    +\int_s^{s+h-v}A'(u+v)\overline{A(u)}\ud u.
  \]
  Thus, $c(h) = 0$, $|c(v)| \le hM^2$ and $|c'(v)| \le M^2+hMD$. Splitting the
  covariance square into its two conjugate triangles gives
  \[
    \|1_IA(t)e^{i\omega t}\|_a^2
    = 2\operatorname{Re}\int_0^h v^{-a}e^{i\omega v}c(v)\ud v
    = 2h^{2-a}\operatorname{Re}\int_0^1 y^{-a}e^{i\omega hy}f(y)\ud y,
  \]
  where $f(y) = c(hy)/h$. Set $K = M+hD$. Then $f(1) = 0$ and $|f|,|f'| \le
  K^2$. If $K = 0$ the assertion is immediate. If $|\omega h| \le 1$, the last
  integral is at most $K^2/(1-a)$ in absolute value, and $\langle\omega
  h\rangle^{a-1} \ge 2^{(a-1)/2}$. The claimed constant is sufficient in this
  case.

  Now suppose $k = |\omega h| > 1$. On $[0,1/k]$ the integral is at most
  $K^2k^{a-1}/(1-a)$. On $[1/k,1]$, integrate the exponential by parts. The
  endpoint at one vanishes, while the other endpoint contributes at most
  $K^2k^{a-1}$. The derivative integral, divided by $k$, is bounded by
  \[
    \frac{K^2}k\left(a\int_{1/k}^1 y^{-a-1}\ud y
    +\int_{1/k}^1 y^{-a}\ud y\right)
    \le K^2 k^{a-1}\left(1+\frac1{1-a}\right).
  \]
  Here $k^{-1} \le k^{a-1}$ since $a > 0$. The total is at most
  $4K^2k^{a-1}/(1-a)$. Finally, $k^{a-1} \le \sqrt{2}\langle\omega
  h\rangle^{a-1}$ and $8\sqrt{2} \le 16$. Including the factor $2h^{2-a}$ proves
  the bound for either sign of $\omega$. This proves
  Lemma~\ref{L:masked-temporal-modulation}.
\end{proof}

\begin{lemma}[Correlated trigonometric products]\label{L:masked-trig-products}
  Let $r,\rho \ge 0$. Let $P_r$ be either $\cos(rt)$ or $\sin(rt)$, and let
  $Q_\rho$ be either $\cos(\rho t)$ or $\sin(\rho t)$. With the amplitude and
  constant above,
  \begin{align}
    \|1_IAP_rQ_\rho\|_a^2
    &\le C_0(a)h^{2-a}(M+hD)^2\mathcal{Q}_h(r,\rho),
    \label{E:masked-trig-products}\\
    \mathcal{Q}_h(r,\rho)
    &:= \langle h(r-\rho)\rangle^{a-1}
    +\langle h(r+\rho)\rangle^{a-1}.
    \label{E:masked-frequency-pair}
  \end{align}
\end{lemma}

\begin{proof}
  Expanding the two trigonometric factors gives four complex exponential modes
  with frequencies $\sigma r+\tau\rho$, $\sigma,\tau \in \{-1,1\}$, and
  coefficients of absolute value $1/4$. In a Hilbert space,
  \[
    \left\|\sum_j c_jv_j\right\|^2
    \le \left(\sum_j|c_j|\right)\sum_j|c_j|\|v_j\|^2.
  \]
  The first coefficient sum is one. Apply
  Lemma~\ref{L:masked-temporal-modulation} to each mode. Two frequencies have
  absolute value $|r-\rho|$ and two have absolute value $r+\rho$. This gives one
  half of the asserted right side, and hence the bound. Repeated or zero
  frequencies are permitted; no orthogonality or independence of the modes is
  used. This proves Lemma~\ref{L:masked-trig-products}.
\end{proof}

\begin{lemma}[Continuous sine factors]\label{L:masked-continuous-wave}
  Define $c_r(t) = \cos(rt)$ and
  \[
    b_r(t) = \langle r\rangle\frac{\sin(rt)}r\quad(r > 0),
    \qquad b_0(t) = t.
  \]
  For $F_r \in \{c_r,b_r\}$ and $G_\rho \in \{c_\rho,b_\rho\}$,
  \begin{equation}\label{E:masked-continuous-wave}
    \|1_IF_rG_\rho\|_a^2
    \le C_1(a)(1+T)^4h^{2-a}\mathcal{Q}_h(r,\rho),
    \qquad C_1(a) = 4\cdot3^{1-a}C_0(a).
  \end{equation}
\end{lemma}

\begin{proof}
  For $r \ge 1$, $b_r$ is a sine times a coefficient at most $\sqrt{2}$. For $0
  \le r \le 1$, the fundamental theorem of calculus gives
  \[
    b_r(t) = t\langle r\rangle\int_0^1\cos(vrt)\ud v.
  \]
  This includes $r = 0$. Hence, a product $F_rG_\rho$ is a scalar coefficient of
  absolute value at most two times a normalized integral of
  $t^jP_{r'}(t)Q_{\rho'}(t)$, where $j \in \{0,1,2\}$ and $|r'-r|,|\rho'-\rho|
  \le 1$. The integral is over zero, one or two copies of $[0,1]$; with no
  averages, use a one-point probability space. Its parameter dependence is
  continuous in $L^2(0,T)$ by dominated convergence.
  Lemma~\ref{L:masked-time-completion} makes it continuous in $H_a$ as well.
  Thus, the averages are Bochner integrals, and their $H_a$ images agree with
  the pointwise averages by continuity of the inclusion. Hilbert Jensen bounds
  the squared norm of an average by the average of the squared norms.

  Apply Lemma~\ref{L:masked-trig-products} to each integrand with $A(t) = t^j$.
  For $j = 0,1,2$, respectively, $M+hD$ is at most $1$, $T+1$ and $T^2+2T$. In
  every case $(M+hD)^2 \le (1+T)^4$. For real $x,e$, the triangle inequality and
  its reverse imply
  \[
    \frac{\langle x\rangle}{1+|e|}
    \le \langle x+e\rangle \le (1+|e|)\langle x\rangle.
  \]
  The changes in $h(r\pm\rho)$ have magnitude at most two, since $h \le 1$.
  Because $a-1 < 0$, each altered bracket to the power $a-1$ is at most
  $3^{1-a}$ times the original one. Apply these comparisons, Hilbert Jensen and
  the squared scalar coefficient bound four to obtain
  \eqref{E:masked-continuous-wave}. This proves
  Lemma~\ref{L:masked-continuous-wave}.
\end{proof}

These estimates concern covariance forms and temporal factors. Their application
to displacement and velocity is in Section~\ref{S:wave-variance}; the subsequent
operator, source and moment results require their separate proofs. They do not
follow from covariance domination alone.
